% =====================================================================
% Chapitre 4 : Génétique des populations vérifiée
% Sources primaires : P7 (doi:10.5281/zenodo.23653850), N4 (doi:10.5281/zenodo.23660071),
% H11/H12 (results/autoresearch/2026-10-25/), Lean : formal_cert/Autoresearch/PopGen*.lean,
% formal_cert/Programmes/PopGen/*.lean, formal_cert/Programmes/EvoGraph/*.lean.
% Les extraits Lean sont recopiés mot pour mot, au rendu des glyphes près.
% =====================================================================
\chapter{Génétique des populations vérifiée : fixation, sélection et graphes}
\label{ch4:debut}\label{ch4:chapitre}
\chaptermark{Génétique des populations vérifiée}
\sloppy

Une mutation apparaît dans une population de $N$ individus. Va-t-elle disparaître, ou finira-t-elle par être
portée par tous ? C'est la question de la \emph{fixation}. Elle est à la fois très concrète (une bactérie
devenue résistante à un antibiotique envahira-t-elle la colonie ?) et parfaitement accessible avec les outils de
la classe préparatoire. Il suffit de connaître les chaînes de Markov finies, les sommes géométriques, la
formule du binôme, les suites récurrentes linéaires et un peu de topologie de $\R$.

Ce chapitre suit le même fil que les précédents. Chaque énoncé reçoit une \emph{démonstration humaine} de niveau
CPGE, puis, à côté, l'énoncé \emph{tel que le noyau de Lean~4 l'a vérifié}, avec le nom de fichier et le numéro
de ligne. Les deux ne coïncident pas toujours : il arrive que la preuve Lean suive un autre chemin que la
preuve humaine, et nous le signalons chaque fois.

\begin{warnbox}[: statut de ce chapitre]
\begin{itemize}[leftmargin=1.2em,itemsep=1pt]
\item Tous les résultats de ce chapitre sont \textbf{classiques} : Moran (1958) \cite{ch4:Mor}, Fisher, Wright,
Price (1970) \cite{ch4:Pri70}, Lieberman--Hauert--Nowak (2005) \cite{ch4:LHN}. Aucune nouveauté mathématique n'est revendiquée. Ce qui est propre au projet,
c'est une formalisation vérifiée par le noyau Lean. Nous n'avons pas fait de recherche systématique de
formalisations antérieures.
\item Les fichiers Lean, les démonstrations et ce texte ont été produits par des agents d'IA (Claude) sous la
direction de l'auteur. Un agent d'IA distinct les a audités : recompilation depuis les sources, audit des
axiomes, contrôles plantés. Les verdicts sont consignés dans \texttt{AUDIT.md} et \texttt{verification.md}.
\textbf{Aucun humain n'a audité ces résultats.}
\item Rien ici ne dépend des théorèmes d'openai/math des chapitres~1 à~3.
\end{itemize}
\end{warnbox}

\paragraph{Rappel des niveaux.} Voici ce que signifie chaque étiquette dans ce chapitre.
\begin{itemize}[itemsep=1pt]
\item \tierA{} signifie que l'\emph{énoncé Lean} cité est vérifié par le noyau. Ses axiomes sont dans
$\{\lean{propext},\lean{Classical.choice},\lean{Quot.sound}\}$, et le module a passé un rejeu noyau
(\texttt{leanchecker}) ainsi qu'une recherche des échappatoires \lean{skipKernelTC}, \lean{addDecl},
\lean{native_decide} et \lean{implemented_by}. Le niveau~A porte sur l'énoncé Lean, \emph{jamais} sur sa
lecture biologique.
\item \tierL{} signifie qu'il s'agit d'une source publiée, lue et localisée à la page près.
\item \tierC{} signifie un argument humain, y compris une démonstration de ce livre, qui n'est pas vérifié par
une machine.
\item \tierX{} signifie un calcul (flottant, ou rationnel exact par un script) qui illustre et contrôle, mais
n'établit jamais rien.
\end{itemize}

\paragraph{Couverture du rejeu noyau.} Les audits d'origine vérifiaient deux choses : la liste des axiomes, et
une recompilation par un agent distinct. L'audit N4 a ajouté un rejeu \texttt{leanchecker} des modules
\lean{Programmes.PopGen.Absorption}, \lean{Programmes.PopGen.MoranModel} et des trois modules
\lean{Programmes.EvoGraph.*}. Son contrôle négatif était un théorème falsifié sous \lean{debug.skipKernelTC},
et le rejeu l'a bien rejeté. Les autres modules cités ici étaient sans rejeu : \lean{Autoresearch.PopGenFisher},
\lean{Autoresearch.PopGenWrightFisher}, \lean{Programmes.PopGen.WrightFisherFixation},
\lean{Programmes.PopGen.MoranFixation} et \lean{Programmes.PopGen.Price}. Pour eux, l'agent rédacteur de ce
chapitre a lancé \texttt{leanchecker} au moment de la rédaction (2026-10-26), sur les \texttt{.olean} existants,
tous plus récents que leurs sources, avec la chaîne épinglée Lean~4.34.1 ; il a aussi cherché par \texttt{grep}
les jetons interdits. Le résultat est donné au §\ref{ch4:sec:bilan}. Ce rejeu n'est pas un audit humain.

\paragraph{Le pont vers les probabilités \tierC.} Dans tout le chapitre, $P$ désigne une matrice stochastique
sur un ensemble fini d'états. On définit $P^n$ par la récurrence de Chapman--Kolmogorov $P^0=I$,
$P^{n+1}(i,j)=\sum_kP^n(i,k)P(k,j)$ (c'est \lean{kPow} en Lean). Si $(X_n)$ est la chaîne de Markov de matrice
$P$ issue de $i$, on a par définition $\mathbb P_i(X_n=j)=P^n(i,j)$. Si de plus l'état $N$ est absorbant,
$\{X_n=N\}=\{T_N\le n\}$, où $T_N$ est le temps d'atteinte de $N$. La limite de $P^n(i,N)$ est donc
$\mathbb P_i(T_N<\infty)$, la \emph{probabilité de fixation}. Cette identification est standard, mais elle
\emph{n'est pas} un énoncé Lean : aucun espace des trajectoires n'est construit. Les théorèmes Lean portent sur
les limites des puissances de matrices, et c'est la formulation que nous retenons, par exemple « la probabilité
d'absorption en $N$ après $n$ pas tend vers \dots ».

% ---------------------------------------------------------------------
\section{Hardy--Weinberg et sélection à un locus}\label{ch4:sec:hw}\sectionmark{Hardy--Weinberg et sélection}

\subsection{Le modèle}
Un gène (un \emph{locus}) possède deux allèles $A_1$ et $A_2$, de fréquences $p$ et $q=1-p$. Les individus sont
diploïdes : chacun porte deux copies, donc l'un des trois génotypes $A_1A_1$, $A_1A_2$ ou $A_2A_2$. Les
générations ne se chevauchent pas. On suppose de plus un accouplement au hasard : un zygote est formé de deux
gamètes tirés indépendamment dans le pool gamétique.

\begin{proposition}[Hardy--Weinberg]\label{ch4:prop:hw}
Partons de fréquences génotypiques $(x_{11},x_{12},x_{22})$ de somme~$1$, et posons $p=x_{11}+x_{12}/2$. Après une
génération d'accouplement au hasard, les fréquences sont $(p^2,2pq,q^2)$, de somme~$1$, et la fréquence de
l'allèle $A_1$ vaut encore $p$. Ces proportions sont stationnaires dès la première génération.
\tierA{} (\lean{hwStep_sum}, \lean{hwStep_allele}, \lean{hwStep_stationary}, \texttt{PopGenFisher.lean},
l.~129--148).
\end{proposition}

\begin{humanproof}
Un gamète porte $A_1$ avec probabilité $x_{11}+\frac12x_{12}=p$. Les deux gamètes étant indépendants, on obtient
$A_1A_1$ avec probabilité $p^2$, $A_1A_2$ avec probabilité $2pq$ (deux ordres possibles) et $A_2A_2$ avec
probabilité $q^2$. La somme vaut $(p+q)^2=1$. Dans la nouvelle génération, la fréquence de $A_1$ est
$p^2+\frac12\cdot2pq=p(p+q)=p$ : l'application est donc idempotente.
\end{humanproof}

\begin{leanbox}[: Hardy--Weinberg (\texttt{formal\_cert/Autoresearch/PopGenFisher.lean}, l.~125--127 et 145)]
\begin{lstlisting}[language=Lean]
noncomputable def hwStep (x : ℝ × ℝ × ℝ) : ℝ × ℝ × ℝ :=
  let p := x.1 + x.2.1 / 2
  (p ^ 2, 2 * p * (1 - p), (1 - p) ^ 2)
theorem hwStep_stationary (x : ℝ × ℝ × ℝ) : hwStep (hwStep x) = hwStep x := by
\end{lstlisting}
L'énoncé \lean{hwStep_stationary} vaut pour tout triplet, même non normalisé, car $p^2+p(1-p)=p$ est une identité.
En revanche, \lean{hwStep} ignore $x_{22}$ : il suppose implicitement l'entrée normalisée (remarque du
vérificateur, \texttt{H12/verification.md}).
\end{leanbox}

\subsection{Sélection de viabilité}
Les génotypes $A_1A_1$, $A_1A_2$ et $A_2A_2$ survivent jusqu'à l'âge adulte avec des \emph{valeurs sélectives}
(fitness) $w_{11},w_{12},w_{22}>0$. Avant sélection, les zygotes sont en proportions de Hardy--Weinberg. On en
déduit la \emph{valeur sélective moyenne} et la nouvelle fréquence allélique :
\[
\bar w(p)=p^2w_{11}+2pq\,w_{12}+q^2w_{22},\qquad
p'=\frac{p^2w_{11}+pq\,w_{12}}{\bar w(p)}.
\]

\begin{leanbox}[: le modèle (\texttt{PopGenFisher.lean}, l.~19--24)]
\begin{lstlisting}[language=Lean]
noncomputable def wbar (w11 w12 w22 p : ℝ) : ℝ :=
  p ^ 2 * w11 + 2 * p * (1 - p) * w12 + (1 - p) ^ 2 * w22

/-- Allele frequency after one generation of viability selection. -/
noncomputable def pnext (w11 w12 w22 p : ℝ) : ℝ :=
  (p ^ 2 * w11 + p * (1 - p) * w12) / wbar w11 w12 w22 p
\end{lstlisting}
\end{leanbox}

\begin{lemma}\label{ch4:lem:dp}
Soient $w_{11},w_{12},w_{22}>0$ et $p\in[0,1]$. Alors :
\begin{enumerate}[label=(\roman*),itemsep=1pt]
\item $\bar w(p)>0$ et $p'\in[0,1]$ ;
\item $\Delta p:=p'-p=\dfrac{pq\,s}{\bar w(p)}$, avec $s=p(w_{11}-w_{12})+q(w_{12}-w_{22})$ ;
\item $\bar w'(p)=2s$, d'où la forme de Wright $\Delta p=\dfrac{pq\,\bar w'(p)}{2\bar w(p)}$.
\end{enumerate}
\tierA{} (\lean{wbar_pos}, \lean{pnext_mem_Icc}, \lean{delta_p_eq}, \lean{hasDerivAt_wbar},
\lean{delta_p_eq_deriv}, \texttt{PopGenFisher.lean}, l.~26--91).
\end{lemma}

\begin{humanproof}
(i) Si $p>0$, le terme $p^2w_{11}$ est $>0$ et les deux autres sont $\ge0$. Si $p=0$, $\bar w=w_{22}>0$. Le
numérateur de $p'$ est $\ge0$, et $\bar w-(p^2w_{11}+pqw_{12})=pqw_{12}+q^2w_{22}\ge0$, donc $p'\le1$.
(ii) On réduit au même dénominateur, en utilisant $1-p=q$ et $1-2p=q-p$ :
\[
p^2w_{11}+pqw_{12}-p\,\bar w
=p\bigl[pw_{11}(1-p)+qw_{12}(1-2p)-q^2w_{22}\bigr]
=pq\bigl[pw_{11}+(q-p)w_{12}-qw_{22}\bigr]=pq\,s ,
\]
puis on divise par $\bar w>0$.
(iii) En développant, $\bar w(x)=w_{22}+2(w_{12}-w_{22})x+(w_{11}-2w_{12}+w_{22})x^2$, donc
$\bar w'(p)=2(w_{12}-w_{22})+2(w_{11}-2w_{12}+w_{22})p=2s$.
\end{humanproof}

En Lean, la dérivée de (iii) n'est pas un polynôme écrit à la main. C'est la dérivée réelle \lean{deriv} de la
fonction $x\mapsto\bar w(x)$, obtenue à partir d'un énoncé \lean{HasDerivAt} explicite (l.~64--82).

\subsection{L'identité exacte de l'incrément}

\begin{theorem}[Monotonie de la valeur sélective moyenne]\label{ch4:thm:wbar}
Soient $w_{11},w_{12},w_{22}>0$ et $p\in[0,1]$. Avec les notations du lemme~\ref{ch4:lem:dp}, on a l'identité
exacte
\begin{equation}\label{ch4:eq:increment}
\bar w(p')-\bar w(p)=\frac{pq\,s^2\,\bigl(\bar w(p)+p\,w_{11}+q\,w_{22}\bigr)}{\bar w(p)^2},
\end{equation}
et par conséquent $\bar w(p')\ge\bar w(p)$ : la valeur sélective moyenne ne décroît jamais d'une génération à la
suivante. \tierA{} (\lean{wbar_increment} et \lean{fisher_mean_fitness_nondecreasing},
\texttt{PopGenFisher.lean}, l.~95--119).
\end{theorem}

\begin{humanproof}
Posons $W=\bar w(p)$, $a=w_{11}-2w_{12}+w_{22}$ et $\Delta=p'-p=pqs/W$ (lemme~\ref{ch4:lem:dp}).

\emph{Étape 1 : Taylor est exact.} $\bar w$ est un polynôme de degré~$2$ en $x$, de dérivée $2s$ en $p$ et de
dérivée seconde $2a$. La formule de Taylor à l'ordre~$2$ est donc exacte :
\[
\bar w(p+\Delta)=W+2s\,\Delta+a\,\Delta^2 .
\]

\emph{Étape 2 : on remplace $\Delta$.}
\[
\bar w(p')-W=2s\cdot\frac{pqs}{W}+a\,\frac{p^2q^2s^2}{W^2}=\frac{pq\,s^2}{W^2}\,\bigl(2W+a\,pq\bigr).
\]

\emph{Étape 3 : l'identité clé.} On vérifie en développant que
\[
W+a\,pq=p^2w_{11}+2pqw_{12}+q^2w_{22}+pq(w_{11}-2w_{12}+w_{22})=pw_{11}(p+q)+qw_{22}(p+q)=pw_{11}+qw_{22}.
\]
Ainsi $2W+a\,pq=W+pw_{11}+qw_{22}$, ce qui donne exactement~\eqref{ch4:eq:increment}.

\emph{Étape 4 : le signe.} Les facteurs $p$, $q$, $s^2$, $W^2$ sont $\ge0$ (et $W>0$), et
$W+pw_{11}+qw_{22}>0$. Le quotient est donc $\ge0$.
\end{humanproof}

\begin{leanbox}[: l'incrément (\texttt{PopGenFisher.lean}, l.~95--99 et 110--112)]
\begin{lstlisting}[language=Lean]
theorem wbar_increment {w11 w12 w22 p : ℝ} (h11 : 0 < w11) (h12 : 0 < w12) (h22 : 0 < w22)
    (hp0 : 0 ≤ p) (hp1 : p ≤ 1) :
    wbar w11 w12 w22 (pnext w11 w12 w22 p) - wbar w11 w12 w22 p
      = p * (1 - p) * (p * (w11 - w12) + (1 - p) * (w12 - w22)) ^ 2
          * (wbar w11 w12 w22 p + p * w11 + (1 - p) * w22) / (wbar w11 w12 w22 p) ^ 2 := by
theorem fisher_mean_fitness_nondecreasing {w11 w12 w22 p : ℝ}
    (h11 : 0 < w11) (h12 : 0 < w12) (h22 : 0 < w22) (hp0 : 0 ≤ p) (hp1 : p ≤ 1) :
    wbar w11 w12 w22 p ≤ wbar w11 w12 w22 (pnext w11 w12 w22 p) := by
\end{lstlisting}
\emph{Côte à côte.} La preuve Lean ne passe pas par Taylor. Elle chasse les dénominateurs (\lean{field_simp}),
puis conclut par \lean{ring}. Le docstring précise que l'identité a d'abord été trouvée avec \texttt{sympy}, puis
vérifiée par le noyau. La preuve humaine ci-dessus explique \emph{pourquoi} le facteur
$\bar w+pw_{11}+qw_{22}$ apparaît : c'est l'identité de l'étape~3.
\end{leanbox}

\begin{remark}[Hypothèses \tierC]\label{ch4:rem:positivite}
La positivité \emph{stricte} des valeurs sélectives est suffisante, mais pas nécessaire : par~\eqref{ch4:eq:increment},
$w\ge0$ avec $\bar w(p)>0$ suffit, puisqu'alors $\bar w+pw_{11}+qw_{22}\ge\bar w>0$. C'est une correction du
vérificateur indépendant (\texttt{H12/verification.md}). Le théorème Lean est correct, mais plus étroit qu'il ne
pourrait l'être. La \emph{non-négativité}, elle, ne peut pas être abandonnée. Le contrôle
\lean{fisher_fails_without_positivity} prend $w=(-1,1,0)$ et $p=\frac12$ : on a $\bar w=\frac14>0$, mais $p'=0$ et
$\bar w(p')=0<\frac14$ \tierA{} (l.~164--167).
\end{remark}

\begin{remark}[La moyenne monte, mais $p$ peut fuir l'équilibre \tierA]
En cas de sous-dominance, $w=(1,\frac12,1)$, il existe un équilibre intérieur $p=\frac12$. Pourtant, en partant de
$p=\frac25$, on obtient $p'=\frac7{19}<\frac25$ : la fréquence \emph{s'éloigne} de l'équilibre, et $\bar w$
augmente quand même ($\frac{19}{25}<\frac{277}{361}$). C'est \lean{underdominance_moves_away} (l.~155--158),
et les valeurs numériques ont été recalculées en rationnels exacts \tierX. La croissance de $\bar w$ ne dit donc
rien de la convergence vers un polymorphisme.
\end{remark}

\begin{warnbox}[: ce n'est pas le « théorème fondamental » de Fisher]
Le théorème~\ref{ch4:thm:wbar} couvre un seul locus, deux allèles, des générations discrètes, une sélection de
viabilité, des valeurs sélectives constantes et l'accouplement au hasard. C'est la monotonie de la valeur
sélective moyenne \emph{pour ce modèle}. Le théorème fondamental de Fisher, au sens de la lecture de Price (1972)
en termes de « changement partiel » \cite{ch4:Pri72} et de variance génétique additive, est un autre énoncé, plus fin, qui
\emph{n'est pas formalisé} ici \tierC. Le docstring Lean d'origine parlait de « Fisher's fundamental theorem ». Le
vérificateur a exigé qu'on le renomme, et l'en-tête du module dit désormais « one-locus, two-allele mean-fitness
monotonicity ».
\end{warnbox}

% ---------------------------------------------------------------------
\section[Wright--Fisher neutre]{Wright--Fisher neutre : martingale, fonctions harmoniques, absorption}\label{ch4:sec:wf}

\subsection{La chaîne}
On s'intéresse maintenant au \emph{hasard} seul. Considérons une population de $N$ gènes, dont $i$ sont de type
mutant. La génération suivante s'obtient en tirant $N$ gènes \emph{avec remise} dans la génération présente.
Sans sélection, le nombre $j$ de mutants dans la génération suivante suit la loi binomiale
$\mathcal B(N,i/N)$ :
\[
P(i,j)=\binom Nj\Bigl(\frac iN\Bigr)^j\Bigl(1-\frac iN\Bigr)^{N-j},\qquad 0\le i,j\le N .
\]

\begin{leanbox}[: le noyau (\texttt{formal\_cert/Autoresearch/PopGenWrightFisher.lean}, l.~25--26)]
\begin{lstlisting}[language=Lean]
noncomputable def wfP (N : ℕ) (i j : Fin (N + 1)) : ℝ :=
  (N.choose (j : ℕ) : ℝ) * ((i : ℕ) / N : ℝ) ^ (j : ℕ) * (1 - ((i : ℕ) / N : ℝ)) ^ (N - (j : ℕ))
\end{lstlisting}
Le contrôle positif \lean{ctl_positive_wfP_value} établit \lean{wfP 2 1 1 = 1 / 2} : la division est bien
réelle, et non la division entière de $\N$.
\end{leanbox}

\begin{proposition}\label{ch4:prop:wfbase}
Soit $N\ge1$.
\begin{enumerate}[label=(\roman*),itemsep=1pt]
\item $P$ est stochastique, et les états $0$ et $N$ sont absorbants.
\item \emph{(Martingale)} Pour tout $i$, $\sum_jP(i,j)\,j=i$. Autrement dit, $h(j)=j/N$ est \emph{harmonique} :
$Ph=h$.
\item Depuis tout état $0\le i<N$, on a $P(i,0)=(1-i/N)^N>0$.
\end{enumerate}
\tierA{} (\lean{wfP_nonneg}, \lean{wfP_sum_one}, \lean{wfP_zero}, \lean{wfP_last}, \lean{wfP_mean},
\lean{frac_harmonic}, \lean{wfP_interior_zero_pos}, \texttt{PopGenWrightFisher.lean}, l.~34--118).
\end{proposition}

\begin{humanproof}
(i) Par la formule du binôme, $\sum_j\binom Nj x^j(1-x)^{N-j}=(x+1-x)^N=1$. En $i=0$ on a $x=0$, et avec la
convention $0^0=1$ la ligne vaut $\delta_0$ ; de même en $i=N$.
(ii) On utilise $j\binom Nj=N\binom{N-1}{j-1}$ pour $j\ge1$. Avec $x=i/N$ et $y=1-x$,
\[
\sum_{j=0}^N\binom Nj x^jy^{N-j}\,j=N\sum_{j=1}^{N}\binom{N-1}{j-1}x^jy^{N-j}=Nx\,(x+y)^{N-1}=Nx=i .
\]
(iii) C'est le terme $j=0$, qui vaut $(1-i/N)^N$, strictement positif dès que $i<N$.
\end{humanproof}

L'identité (ii) dit que le nombre $X_n$ de mutants est une \emph{martingale} : son espérance conditionnelle à
l'étape suivante est sa valeur présente \tierC{} (lecture probabiliste). Par récurrence sur $n$, on obtient
$\sum_jP^n(i,j)\,j=i$ pour tout $n$. Ce calcul ne demande aucune théorie de la mesure.

\subsection{Fonctions harmoniques et principe du maximum}
\begin{definition}
Soit $P$ une matrice stochastique sur $\{0,\dots,N\}$ dont les états $0$ et $N$ sont absorbants. Une fonction
$g:\{0,\dots,N\}\to\R$ est \emph{harmonique à l'intérieur} si $\sum_jP(i,j)g(j)=g(i)$ pour tout $0<i<N$.
\end{definition}

L'intérêt de cette notion est le suivant. La probabilité de fixation $u(i)=\mathbb P_i(T_N<\infty)$ vérifie
\emph{formellement} l'équation « au premier pas » $u(i)=\sum_jP(i,j)u(j)$, avec $u(0)=0$ et $u(N)=1$ \tierC.
Si l'on sait qu'une seule fonction harmonique a ces valeurs au bord, on connaît donc $u$.

\begin{theorem}[Principe du maximum, unicité]\label{ch4:thm:maxprinc}
Soit $N\ge1$, et soit $d$ une fonction harmonique à l'intérieur pour le noyau de Wright--Fisher, avec $d(0)=d(N)=0$.
Alors $d\le0$, et en appliquant le même raisonnement à $-d$, $d=0$. Par conséquent, $i\mapsto i/N$ est
\emph{l'unique} fonction harmonique à l'intérieur qui vaut $0$ en $0$ et $1$ en $N$. \tierA{} (\lean{max_principle},
\lean{harmonic_unique}, \lean{wright_fisher_fixation_harmonic}, \texttt{PopGenWrightFisher.lean}, l.~122--187).
\end{theorem}

\begin{humanproof}
Soit $m$ un point où $d$ atteint son maximum ; il existe, car l'ensemble est fini. Si $m\in\{0,N\}$, alors
$\max d=0$. Sinon $m$ est intérieur, et
\[
0=d(m)-\sum_jP(m,j)d(j)=\sum_jP(m,j)\bigl(d(m)-d(j)\bigr),
\]
où l'on a utilisé $\sum_jP(m,j)=1$. Chaque terme est $\ge0$, puisque $P\ge0$ et que $d(m)$ est maximal ; ils
sont donc tous nuls. Or $P(m,0)>0$ (proposition~\ref{ch4:prop:wfbase}(iii)), donc $d(m)=d(0)=0$. Dans les deux
cas, $\max d=0$. Pour l'unicité, si $g$ convient, $d=g-h$ avec $h(j)=j/N$ est harmonique à l'intérieur et
s'annule au bord.
\end{humanproof}

\begin{leanbox}[: le principe du maximum (\texttt{PopGenWrightFisher.lean}, l.~122--126)]
\begin{lstlisting}[language=Lean]
theorem max_principle {N : ℕ} (hN : 0 < N) (d : Fin (N + 1) → ℝ) (h0 : d 0 = 0)
    (hN' : d (Fin.last N) = 0)
    (hharm : ∀ i : Fin (N + 1), 0 < (i : ℕ) → (i : ℕ) < N → d i = ∑ j, wfP N i j * d j) :
    ∀ i, d i ≤ 0 := by
\end{lstlisting}
La preuve Lean suit exactement la preuve humaine. \lean{Finset.exists_max_image} fournit le point de maximum, puis
\lean{Finset.sum_eq_zero_iff_of_nonneg} annule chaque terme, et \lean{wfP_interior_zero_pos} conclut.
\end{leanbox}

\begin{warnbox}[: ce que H11 ne prouvait pas]
Le module H11 \cite{ch4:H1112} s'arrête là. Il établit que $i/N$ est \emph{la} solution harmonique, mais le pas « la probabilité
de fixation est harmonique » n'y est pas formalisé. Le vérificateur l'a écrit sans détour : « this file does NOT
kernel-check ``fixation probability $= i/N$'' » (\texttt{H11/NOTE.md}). Il proposait de combler la lacune
\emph{sans} théorie de la mesure, en démontrant $P^n(i,N)\to i/N$. C'est ce qu'a fait le programme P7, et c'est
l'objet du paragraphe suivant.
\end{warnbox}

\subsection{L'absorption : \texorpdfstring{$P^n(i,N)\to i/N$}{P\^{}n(i,N) -> i/N}}
On note $I_n(i)=\sum_{0<k<N}P^n(i,k)$ la \emph{masse intérieure} après $n$ pas, c'est-à-dire la probabilité de ne
pas encore être absorbé (\lean{interiorMass}).

\begin{theorem}[Fixation neutre]\label{ch4:thm:wflim}
Pour tout $N\ge1$ et tout $0\le i\le N$,
\[
P^n(i,N)\xrightarrow[n\to\infty]{}\frac iN,\qquad P^n(i,0)\xrightarrow[n\to\infty]{}1-\frac iN,
\]
et la suite $\bigl(P^n(i,N)\bigr)_n$ est croissante. \tierA{} (\lean{wright_fisher_fixation_limit},
\lean{wright_fisher_fixation_mono}, \texttt{Programmes/PopGen/}\allowbreak\texttt{WrightFisherFixation.lean}, l.~49--58). Le théorème général \lean{fixation_limit} dont il découle reparaît sous le nom \lean{KAL}.
\end{theorem}

Nous donnons d'abord une démonstration de niveau CPGE \tierC, qui repose sur la \emph{décroissance géométrique}
de la masse intérieure. La preuve Lean suit un autre chemin, présenté ensuite.

\begin{lemma}[Décroissance géométrique \tierC]\label{ch4:lem:geom}
Posons $\delta=\min_{0<k<N}P(k,0)$. Alors $\delta\ge N^{-N}>0$ et, pour tout $n$, $I_{n+1}(i)\le(1-\delta)\,I_n(i)$.
Par conséquent $I_n(i)\le(1-\delta)^n$.
\end{lemma}

\begin{humanproof}
Pour $0<k<N$, on a $P(k,0)=(1-k/N)^N\ge(1/N)^N$. Ensuite,
\[
I_{n+1}(i)=\sum_{0<j<N}\sum_kP^n(i,k)P(k,j)=\sum_kP^n(i,k)\sum_{0<j<N}P(k,j).
\]
Si $k\in\{0,N\}$, la somme intérieure est nulle, car l'état est absorbant. Si $k$ est intérieur, elle vaut au plus
$1-P(k,0)\le1-\delta$. Donc $I_{n+1}(i)\le(1-\delta)\sum_{0<k<N}P^n(i,k)=(1-\delta)I_n(i)$, puis on conclut par
récurrence, puisque $I_0\le1$.
\end{humanproof}

\begin{humanproof}[Démonstration humaine du théorème~\ref{ch4:thm:wflim}]
Posons $h(j)=j/N$. Par la proposition~\ref{ch4:prop:wfbase}, $Ph=h$ en tout état, donc $P^nh=h$, ce qui s'écrit
\[
\frac iN=\sum_jP^n(i,j)\,\frac jN=P^n(i,N)+\sum_{0<j<N}P^n(i,j)\,\frac jN .
\]
Le dernier terme est compris entre $0$ et $I_n(i)$, qui tend vers $0$ par le lemme~\ref{ch4:lem:geom}. On en tire
$P^n(i,N)\to i/N$ par encadrement. Ensuite, $P^n(i,0)+P^n(i,N)+I_n(i)=1$ (les lignes de $P^n$ sont de somme~$1$),
d'où $P^n(i,0)\to1-i/N$. Enfin, $P^{n+1}(i,N)=\sum_kP^n(i,k)P(k,N)\ge P^n(i,N)P(N,N)=P^n(i,N)$.
\end{humanproof}

\begin{leanbox}[: le théorème général et le cas Wright--Fisher]
\texttt{formal\_cert/Programmes/PopGen/Absorption.lean}, l.~32--34, 45--53 et 302--308 :
\begin{lstlisting}[language=Lean]
noncomputable def kPow (P : Fin (N + 1) → Fin (N + 1) → ℝ) : ℕ → Fin (N + 1) → Fin (N + 1) → ℝ
  | 0 => fun i j => if i = j then 1 else 0
  | n + 1 => fun i j => ∑ k, kPow P n i k * P k j
structure AbsorbingChain (P : Fin (N + 1) → Fin (N + 1) → ℝ) : Prop where
  nonneg : ∀ i j, 0 ≤ P i j
  rowSum : ∀ i, ∑ j, P i j = 1
  absorb0 : ∀ j, P 0 j = if j = 0 then 1 else 0
  absorbN : ∀ j, P (Fin.last N) j = if j = Fin.last N then 1 else 0
def InteriorHarmonic (P : Fin (N + 1) → Fin (N + 1) → ℝ) (g : Fin (N + 1) → ℝ) : Prop :=
  ∀ i : Fin (N + 1), 0 < (i : ℕ) → (i : ℕ) < N → ∑ j, P i j * g j = g i
theorem fixation_limit {P : Fin (N + 1) → Fin (N + 1) → ℝ} (hP : AbsorbingChain P) (hN : 0 < N)
    {h : Fin (N + 1) → ℝ} (hh : InteriorHarmonic P h) (h0 : h 0 = 0) (h1 : h (Fin.last N) = 1)
    (hvar : ∀ k : Fin (N + 1), 0 < (k : ℕ) → (k : ℕ) < N → 0 < ∑ j, P k j * (h j - h k) ^ 2)
    (i : Fin (N + 1)) :
    Tendsto (fun n => kPow P n i (Fin.last N)) atTop (𝓝 (h i)) ∧
    Tendsto (fun n => kPow P n i 0) atTop (𝓝 (1 - h i)) := by
\end{lstlisting}
\texttt{WrightFisherFixation.lean}, l.~49--51 :
\begin{lstlisting}[language=Lean]
theorem wright_fisher_fixation_limit {N : ℕ} (hN : 0 < N) (i : Fin (N + 1)) :
    Tendsto (fun n => kPow (wfP N) n i (Fin.last N)) atTop (𝓝 (((i : ℕ) : ℝ) / N)) ∧
    Tendsto (fun n => kPow (wfP N) n i 0) atTop (𝓝 (1 - ((i : ℕ) : ℝ) / N)) := by
\end{lstlisting}
\emph{Côte à côte : deux chemins différents.} La preuve Lean n'utilise pas le lemme~\ref{ch4:lem:geom}. Elle
emploie un argument de \emph{second moment}, valable pour toute chaîne absorbante pourvu qu'une fonction
harmonique $h$ ait une variance à un pas $v(k)=\sum_jP(k,j)(h(j)-h(k))^2$ strictement positive en chaque état
intérieur (hypothèse \lean{hvar}). Voici les étapes.
\begin{enumerate}[itemsep=1pt]
\item On a l'identité $\sum_jP(k,j)h(j)^2=h(k)^2+v(k)$ (\lean{second_moment}, l.~145).
\item Avec $\varepsilon=\min v>0$, on obtient $P^{m+1}(h^2)\ge P^m(h^2)+\varepsilon I_m$. En sommant,
$\varepsilon\sum_{m<n}I_m(i)\le\sum_jh(j)^2$.
\item La série $\sum I_m$ est donc à termes positifs et de sommes partielles bornées : elle converge, et
$I_m\to0$ (\lean{interiorMass_tendsto_zero}, l.~156).
\item On conclut par l'encadrement ci-dessus (\lean{absorption_tendsto}, \lean{mass_split}).
\end{enumerate}
Ce chemin ne donne pas la vitesse géométrique. En revanche, il s'applique tel quel au modèle de Moran
(§\ref{ch4:sec:moran}) et aux graphes (§\ref{ch4:sec:graphes}), où le saut direct vers $0$ est impossible. Pour
Wright--Fisher, la variance est positive grâce au seul saut vers~$0$ : $v(k)\ge P(k,0)(k/N)^2>0$
(\lean{wf_var_pos}, l.~34).
\end{leanbox}

\begin{remark}[Contrôles \tierA]\label{ch4:rem:wfctl}
\emph{Contrôle positif.} On a $P^2(1,2)=\frac38$ pour $N=2$ (\lean{ctl_pos_wf_two_steps}, l.~62). La limite
$\frac12$ n'est donc pas encore atteinte au bout de deux générations.
\emph{Contrôle négatif.} L'hypothèse de variance ne peut pas être supprimée. La matrice identité sur $\{0,1,2\}$
est stochastique et absorbante, et $j/2$ est harmonique pour elle ; pourtant $P^n(1,2)=0$ pour tout $n$
(\lean{ctl_neg_identity_no_absorption}, \texttt{Absorption.lean}, l.~363).
\emph{Unicité.} Si $g$ est une autre fonction harmonique de mêmes valeurs au bord, le même argument donne
$P^n(i,N)\to g(i)$. Par unicité de la limite, $g=h$ (\lean{harmonic_unique_of_var}, l.~320). Cela remplace le
principe du maximum par l'unicité des limites.
\end{remark}

\begin{remark}[Le dépôt Genetic]\label{ch4:rem:genetic}
Le dépôt compagnon \texttt{SocrateAI-Scientific-Genetic} \cite{ch4:Genetic} contient une version antérieure de ces résultats, sur
$\mathbb Q$. Elle figure dans \texttt{formal/ANSE/}\allowbreak\texttt{AlixPopulationGenetics.lean} : le principe du maximum
\lean{ilene_harmonic_le_zero} (l.~82), l'unicité \lean{xav_wf_harmonic_solution_unique} (l.~199) et la majoration
sans passage à la limite \lean{xav_wf_fixation_by_time_le} (l.~268), qui affirme $P^t(i,N)\le i/N$ pour tout $t$.
Cette dernière s'obtient en une ligne à partir de la martingale :
$P^t(i,N)\cdot1\le\sum_jP^t(i,j)\,j/N=i/N$. La preuve de cette majoration en une ligne est une démonstration humaine \tierC. Le fichier
n'a pas été recompilé pour ce livre, et nous ne lui attribuons aucun niveau. Les énoncés de ce chapitre
proviennent des modules H11 et P7 d'AutoevolveAI. Entre les deux dépôts, le lien n'est que textuel : aucun terme
Lean ne relie les deux espaces de noms \tierC.
\end{remark}

% ---------------------------------------------------------------------
\section{Le processus de Moran avec sélection}\label{ch4:sec:moran}\sectionmark{Moran avec sélection}

\subsection{Le modèle}
Une population compte $N$ individus, dont $i$ mutants de valeur sélective $r>0$ ; les résidents ont la valeur
sélective $1$. À chaque pas de temps, deux tirages indépendants ont lieu.
\begin{itemize}[itemsep=1pt]
\item Un individu est choisi pour se reproduire, avec une probabilité proportionnelle à sa valeur sélective. Un
mutant est donc choisi avec probabilité $\frac{ri}{ri+N-i}$.
\item Un individu est choisi \emph{uniformément} parmi les $N$ pour mourir, et le descendant prend sa place.
\end{itemize}
Le nombre de mutants monte de $1$ avec probabilité $p_i$, descend de $1$ avec probabilité $q_i$ et reste
inchangé sinon :
\[
p_i=\frac{ri}{ri+N-i}\cdot\frac{N-i}{N},\qquad q_i=\frac{N-i}{ri+N-i}\cdot\frac iN .
\]
Ce sont les taux de Keller et Uğurlu, à une réécriture d'une ligne près \tierL{}~\cite{ch4:KU} (p.~2--4) \tierC.

\begin{leanbox}[: le modèle de Moran (\texttt{formal\_cert/Programmes/PopGen/MoranModel.lean}, l.~30--32, 74--83, 152--155)]
\begin{lstlisting}[language=Lean]
noncomputable def bdP (N : ℕ) (p q : ℕ → ℝ) (i j : Fin (N + 1)) : ℝ :=
  (if (j : ℕ) = (i : ℕ) + 1 then p i else 0) + (if (j : ℕ) + 1 = (i : ℕ) then q i else 0) +
    (if (j : ℕ) = (i : ℕ) then 1 - p i - q i else 0)
noncomputable def moranUp (N : ℕ) (r : ℝ) (i : ℕ) : ℝ :=
  r * i / (r * i + ((N : ℝ) - i)) * (((N : ℝ) - i) / N)
noncomputable def moranDown (N : ℕ) (r : ℝ) (i : ℕ) : ℝ :=
  ((N : ℝ) - i) / (r * i + ((N : ℝ) - i)) * ((i : ℝ) / N)
noncomputable def moranP (N : ℕ) (r : ℝ) : Fin (N + 1) → Fin (N + 1) → ℝ :=
  bdP N (moranUp N r) (moranDown N r)
noncomputable def geomS (r : ℝ) (n : ℕ) : ℝ := ∑ k ∈ range n, r⁻¹ ^ k
noncomputable def moranFix (N : ℕ) (r : ℝ) (i : ℕ) : ℝ := geomS r i / geomS r N
\end{lstlisting}
L'auditeur indépendant a vérifié symboliquement que \lean{moranUp} et \lean{moranDown} coïncident avec les taux
de Keller--Uğurlu (PG-C-0004). C'est un jugement d'IA, et non un audit humain.
\end{leanbox}

\begin{lemma}\label{ch4:lem:moranbase}
Soient $N\ge1$ et $r>0$.
\begin{enumerate}[label=(\roman*),itemsep=1pt]
\item $p_i=r\,q_i$ pour tout $i$ : le rapport des taux est la valeur sélective.
\item $1-p_i-q_i=\dfrac{ri^2+(N-i)^2}{(ri+N-i)N}\ge0$, donc la matrice est stochastique. De plus
$p_0=q_0=p_N=q_N=0$, donc $0$ et $N$ sont absorbants.
\item Pour $0<i<N$, on a $p_i>0$ et $q_i>0$.
\end{enumerate}
\tierA{} (\lean{moranUp_eq}, \lean{moran_stay_eq}, \lean{moran_stay_nonneg}, \lean{moranUp_pos},
\lean{moranDown_pos}, \lean{moran_absorbing}, \texttt{MoranModel.lean}, l.~85--149, et
\texttt{MoranFixation.lean}, l.~25).
\end{lemma}

\begin{humanproof}
(i) On a $p_i/q_i=\frac{ri(N-i)}{(N-i)i}=r$ si $0<i<N$, et les deux taux sont nuls aux bords.
(ii) On calcule
$(ri+N-i)N-ri(N-i)-(N-i)i=riN-ri(N-i)+(N-i)(N-i)=ri^2+(N-i)^2$.
(iii) Tous les facteurs sont strictement positifs quand $0<i<N$.
\end{humanproof}

\subsection{La ruine du joueur}

\begin{theorem}[Fixation dans le modèle de Moran]\label{ch4:thm:moran}
Soient $N\ge1$, $r>0$, et $S_n=\sum_{k<n}r^{-k}$. Pour tout $0\le i\le N$,
\[
P^n(i,N)\xrightarrow[n\to\infty]{}\frac{S_i}{S_N}=
\begin{cases}\dfrac{1-r^{-i}}{1-r^{-N}}&\text{si }r\ne1,\\[2mm] \dfrac iN&\text{si }r=1,\end{cases}
\qquad P^n(i,0)\longrightarrow1-\frac{S_i}{S_N}.
\]
De plus, $S_i/S_N$ est l'unique fonction harmonique à l'intérieur qui vaut $0$ en $0$ et $1$ en $N$, et
$\frac{1-r^{-i}}{1-r^{-N}}\to\frac iN$ quand $r\to1$, $r\ne1$. \tierA{} (\lean{moran_fixation_limit},
\lean{moran_fixation_limit_closed}, \lean{moran_fixation_limit_neutral}, \lean{moran_fixation_unique},
\texttt{MoranFixation.lean}, l.~72--97 ; \lean{moranFix_closed}, \lean{moran_closed_tendsto_neutral},
\texttt{MoranModel.lean}, l.~203 et 231).
\end{theorem}

\begin{humanproof}
\emph{Étape 1 : la récurrence.} Soit $h$ harmonique à l'intérieur, avec $h(0)=0$ et $h(N)=1$. Pour $0<i<N$,
l'équation au premier pas s'écrit
\[
p_ih(i+1)+q_ih(i-1)+(1-p_i-q_i)h(i)=h(i)\iff p_i\bigl(h(i+1)-h(i)\bigr)=q_i\bigl(h(i)-h(i-1)\bigr).
\]
Comme $p_i=rq_i$ et $q_i>0$, les accroissements $\Delta_i=h(i+1)-h(i)$ vérifient $\Delta_i=r^{-1}\Delta_{i-1}$.
C'est une suite géométrique : $\Delta_i=r^{-i}\Delta_0$. En sommant, $h(i)=\Delta_0\,S_i$. La condition $h(N)=1$
impose alors $\Delta_0=1/S_N$, avec $S_N>0$. Il y a donc \emph{au plus} une solution, et réciproquement
$h(i)=S_i/S_N$ en est une, puisque les équivalences se remontent. C'est l'argument classique de la « ruine du
joueur ».

\emph{Étape 2 : forme close.} Si $r\ne1$, posons $\rho=1/r\ne1$. Alors $S_n=\frac{1-\rho^n}{1-\rho}$, d'où
$\frac{S_i}{S_N}=\frac{1-r^{-i}}{1-r^{-N}}$. Si $r=1$, $S_n=n$. Chaque $S_n(r)$ est une somme finie de fonctions
continues de $r$ et $S_N(1)=N\ne0$, donc le quotient est continu en $r=1$. C'est la limite neutre.

\emph{Étape 3 : l'absorption.} On reprend la preuve du théorème~\ref{ch4:thm:wflim} avec $h=S_\cdot/S_N$ :
$P^nh=h$, puis $|h(i)-P^n(i,N)|\le I_n(i)$, car $0\le h\le1$. Il reste à voir que $I_n(i)\to0$. Le
lemme~\ref{ch4:lem:geom} ne s'applique plus tel quel : depuis $i$, on ne peut pas sauter en $0$ en un pas si
$i\ge2$. On raisonne donc par blocs. Depuis tout état intérieur $k$, le chemin $k\to k-1\to\dots\to0$ a une
probabilité $q_kq_{k-1}\cdots q_1>0$ et ne dure que $k\le N-1$ pas. Une fois en $0$, on y reste. Posons
$m=N-1$ et $\eta=\min_{0<k<N}q_k\cdots q_1>0$. On a alors $P^m(k,0)\ge\eta$ pour tout $k$ intérieur, et le même
calcul que dans le lemme~\ref{ch4:lem:geom}, appliqué à la matrice $P^m$, donne $I_{n+m}(i)\le(1-\eta)I_n(i)$.
Comme $I_n$ est décroissante (une masse absorbée le reste), $I_n(i)\to0$ \tierC.
\end{humanproof}

\begin{leanbox}[: le théorème et sa forme close (\texttt{MoranFixation.lean}, l.~72--74 ; \texttt{MoranModel.lean}, l.~173--175, 203--204)]
\begin{lstlisting}[language=Lean]
theorem moran_fixation_limit {N : ℕ} {r : ℝ} (hN : 0 < N) (hr : 0 < r) (i : Fin (N + 1)) :
    Tendsto (fun n => kPow (moranP N r) n i (Fin.last N)) atTop (𝓝 (moranFix N r i)) ∧
    Tendsto (fun n => kPow (moranP N r) n i 0) atTop (𝓝 (1 - moranFix N r i)) :=
theorem moranFix_harmonic_nat {N : ℕ} {r : ℝ} (hr : 0 < r) {i : ℕ} (h1 : 0 < i) :
    moranUp N r i * moranFix N r (i + 1) + moranDown N r i * moranFix N r (i - 1)
      + (1 - moranUp N r i - moranDown N r i) * moranFix N r i = moranFix N r i := by
theorem moranFix_closed {N : ℕ} {r : ℝ} (hr1 : r ≠ 1) (i : ℕ) :
    moranFix N r i = (1 - r⁻¹ ^ i) / (1 - r⁻¹ ^ N) := by
\end{lstlisting}
\emph{Côte à côte.} Lean ne \emph{résout} pas la récurrence. Il \emph{vérifie} que la solution devinée
$S_i/S_N$ est harmonique (\lean{moranFix_harmonic_nat}, en une \lean{linear_combination}). Puis il applique le
théorème général \lean{fixation_limit} du §\ref{ch4:sec:wf}, avec la variance
$v(i)\ge p_i\,(r^{-i}/S_N)^2>0$ (\lean{moran_var_pos_nat}, l.~189). L'argument par blocs de l'étape~3 n'est pas
formalisé ; le second moment le remplace. L'unicité vient de l'unicité des limites
(\lean{moran_fixation_unique}), non de l'étape~1.
\end{leanbox}

\begin{remark}[Contrôles \tierA]\label{ch4:rem:moranctl}
Les contrôles suivants sont tous vérifiés par le noyau.
\begin{itemize}[itemsep=1pt]
\item \emph{Contrôles positifs.} Pour $N=2$, $r=2$, on a $p_1=\frac13$, $q_1=\frac16$ et $h(1)=\frac23$
(\lean{ctl_pos_moran_rates}, \lean{ctl_pos_moranFix_N2}). Pour $N=3$, $r=2$, on a $h=(0,\frac47,\frac67,1)$, et
l'équation au premier pas est vérifiée (\lean{ctl_pos_moranFix_N3}, \texttt{MoranModel.lean}, l.~240--253).
\item \emph{Contrôles négatifs.} La formule « aux exposants inversés » $(1-r^i)/(1-r^N)$ donne $\frac17\ne\frac47$
et viole l'équation au premier pas. La valeur neutre $i/N$ la viole aussi sous sélection (l.~259--276). Au niveau
de la chaîne, $P^n(1,3)$ pour $N=3$, $r=2$ ne tend ni vers $\frac13$ ni vers $\frac17$
(\lean{ctl_neg_wrong_limits}, \texttt{MoranFixation.lean}, l.~102).
\item \emph{Contrôles plantés.} Dans un fichier séparé que rien n'importe, un « théorème » prouvé par \lean{sorry}
et un autre tiré d'un axiome planté affirment les limites fausses $\frac17$ et $\frac13$. Ce fichier
(\texttt{PopGenProgControls.lean}) compile avec le code de retour~$0$, mais \texttt{\#print axioms} affiche \lean{sorryAx}, resp.\ l'axiome planté
\lean{ctl_cheat}. Ils sont donc détectés (PG-A-0012).
\end{itemize}
\end{remark}

\begin{numbox}[: CVODE et le certificat « fail-closed »]
\tierX{} Le programme P7 a aussi intégré l'équation de Kolmogorov uniformisée $\dot\pi=\pi(P-I)$ avec CVODE (BDF, jacobien
exact), dans un nouveau crate \texttt{qf-popgen} de rusty-SUNDIALS. Le crate s'entoure d'un certificat (« fail-closed » : qui refuse par défaut) qui refuse
tout résultat douteux : noyau valide, conservation de la masse, positivité composante par composante,
décroissance de la masse intérieure et accord avec la cible. Bilan : $106$ exécutions acceptées s'accordent avec
les formes closes à $3{,}1\times10^{-10}$ près ; les $9$ contrôles négatifs sont tous refusés. La loi
uniformisée est un mélange de Poisson des $P^n$ \tierC. Ces calculs testent donc les mêmes formules, mais pas les
puissances de matrices Lean, et \emph{ils n'établissent rien}.
\end{numbox}

\subsection{Interprétation : une mutation de résistance envahit-elle la population ?}

\begin{physbox}[: lecture biologique]
\tierC{} On lit souvent le modèle de Moran comme un modèle de la question suivante : une bactérie ou une cellule devient
résistante à un médicament, et sa lignée prend-elle le dessus ? Si la mutation donne, en présence du
médicament, un avantage relatif $r>1$, la probabilité qu'un mutant unique se fixe vaut
\[
\rho_N(1)=\frac{1-r^{-1}}{1-r^{-N}}\xrightarrow[N\to\infty]{}1-\frac1r .
\]
La forme close est \tierA{} (théorème~\ref{ch4:thm:moran}) et la limite en $N$ est un calcul élémentaire \tierC.
Évaluée en flottants \tierX, elle donne pour $N=100$ :
\begin{center}\small
\begin{tabular}{lccc}
\toprule
& $r=0{,}99$ & $r=1{,}01$ & $r=1{,}1$\\
\midrule
$\rho_{100}(1)$ & $0{,}0058$ & $0{,}0157$ & $0{,}0909$\\
valeur neutre $1/N$ & $0{,}01$ & $0{,}01$ & $0{,}01$\\
\bottomrule
\end{tabular}
\end{center}
Deux leçons qualitatives en découlent. D'abord, même un mutant avantagé ($r=1{,}1$) disparaît le plus souvent, ici
dans environ $91\,\%$ des cas, emporté par la dérive. Ensuite, un mutant désavantagé ($r<1$) a encore une petite
chance de se fixer.

\emph{Prudence.} Ce n'est qu'une lecture, et elle reste au niveau~C. Le modèle suppose une population de taille
fixe et bien mélangée, une valeur sélective constante, aucune nouvelle mutation, la mort uniforme et des
générations qui se chevauchent pas à pas. Une infection réelle a une taille qui varie, des concentrations de
médicament qui changent, un coût de la résistance, une structure spatiale et des mutations récurrentes. Aucune
donnée n'a été utilisée, et aucun théorème Lean ne certifie qu'une population réelle suit cette chaîne.
\end{physbox}

% ---------------------------------------------------------------------
\section{L'équation de Price et le corollaire haploïde}\label{ch4:sec:price}\sectionmark{L'équation de Price}

\subsection{Une identité comptable}
Une population finie est formée de types $i\in s$. Le type $i$ a un effectif $n_i$, une valeur sélective $w_i$
(nombre moyen de descendants par individu), un caractère parental $z_i$, et ses descendants ont un caractère moyen
$z'_i$. On note
\[
M=\sum n_i,\quad A=\sum n_iw_i,\quad B=\sum n_iz_i,\quad C=\sum n_iw_iz_i,\quad D=\sum n_iw_iz'_i .
\]
On définit alors $\bar w=A/M$ et $\bar z=B/M$. La moyenne de la génération suivante est $\bar z'=D/A$, puisque le
type $i$ laisse $n_iw_i$ descendants : c'est l'\emph{hypothèse de modélisation}. On note enfin
$\mathrm{Cov}_n(w,z)=\frac1M\sum n_i(w_i-\bar w)(z_i-\bar z)$ et $\mathbb E_n[f]=\frac1M\sum n_if_i$.

\begin{theorem}[Équation de Price]\label{ch4:thm:price}
Sur un corps quelconque, si $M\ne0$ et $A\ne0$, on a
\[
\bar w\,(\bar z'-\bar z)=\mathrm{Cov}_n(w,z)+\mathbb E_n\bigl[w\,(z'-z)\bigr].
\]
\tierA{} (\lean{price_equation}, \texttt{formal\_cert/Programmes/PopGen/Price.lean}, l.~95--105 ; version réelle
sous hypothèses de positivité \lean{price_equation_pos}, l.~170).
\end{theorem}

\begin{humanproof}
D'une part, $\mathrm{Cov}_n(w,z)=\frac1M\bigl(C-\bar wB-\bar zA+M\bar w\bar z\bigr)=\frac{C-AB/M}{M}$ ; d'autre
part, $\mathbb E_n[w(z'-z)]=\frac{D-C}{M}$. La somme vaut $\frac DM-\frac{AB}{M^2}$. Or
\[
\bar w(\bar z'-\bar z)=\frac AM\Bigl(\frac DA-\frac BM\Bigr)=\frac DM-\frac{AB}{M^2}. \qedhere
\]
\end{humanproof}

\begin{leanbox}[: l'énoncé (\texttt{Price.lean}, l.~95--98)]
\begin{lstlisting}[language=Lean,escapeinside={(*@}{@*)}]
theorem price_equation (s : Finset (*@$\iota$@*)) (n w z z' : (*@$\iota$@*) → K)
    (hN : ∑ i ∈ s, n i ≠ 0) (hW : ∑ i ∈ s, n i * w i ≠ 0) :
    wmean s n w * (zbarNext s n w z' - wmean s n z)
      = wcov s n w z + wmean s n (fun i => w i * (z' i - z i)) := by
\end{lstlisting}
La preuve Lean suit la preuve humaine. Elle développe les deux numérateurs en les cinq sommes brutes
(\lean{cov_num_expand}, \lean{trans_num_expand}), puis conclut par un lemme scalaire où $M,A,B,C,D$ sont des
atomes (\lean{price_scalar} : \lean{field_simp} puis \lean{ring}). L'énoncé vaut sur un corps \lean{K} quelconque,
par exemple $\mathbb Q$ ou $\R$, et pour tout ensemble fini de types.
\end{leanbox}

\subsection{Le corollaire haploïde}
\begin{corollary}[Incrément de sélection haploïde]\label{ch4:cor:haploide}
Prenons pour caractère la valeur sélective elle-même, $z=w$, et supposons qu'il n'y a \emph{aucun biais de
transmission}, $z'=w$ : types clonaux, valeurs sélectives constantes, pas de mutation. Alors
\[
\bar w\,(\bar w'-\bar w)=\mathrm{Var}_n(w),\qquad\text{donc}\qquad \Delta\bar w=\frac{\mathrm{Var}_n(w)}{\bar w}\ge0
\quad(\text{si }n_i>0,\ w_i>0).
\]
\tierA{} (\lean{fisher_haploid}, \lean{fisher_haploid_div}, \lean{fisher_haploid_nondecreasing},
\lean{fisher_haploid_pos}, \texttt{Price.lean}, l.~125--185 ; \lean{fisher_haploid_nondecreasing} est ré-exporté sous le nom \lean{AURELIA}).
\end{corollary}

\begin{humanproof}
Avec $z=z'=w$, le terme de transmission $\mathbb E_n[w(w-w)]$ est identiquement nul, et
$\mathrm{Cov}_n(w,w)=\mathrm{Var}_n(w)$. On divise par $\bar w\ne0$. Le signe vient de $\mathrm{Var}_n(w)\ge0$ et de
$\bar w>0$.
\end{humanproof}

C'est l'équation~(2.27) de LaGatta, avec un terme environnemental nul \tierL{}~\cite{ch4:LaG}. C'est aussi
l'analogue discret de la remarque de Baez : $d\bar f/dt=\mathrm{Var}(f)$ lorsque les valeurs sélectives sont
constantes \tierL{}~\cite{ch4:Bae}.

Pour deux allèles de fréquences $p$, $1-p$ et de valeurs sélectives $w_1$, $w_2$, on trouve
$\mathrm{Var}_n(w)=p(1-p)(w_1-w_2)^2$, d'où $\Delta\bar w=p(1-p)(w_1-w_2)^2/\bar w$. Le module P7 en tire, au
cas \lean{Fin 2}, deux énoncés au texte identique à deux lemmes du dépôt Genetic
(\texttt{Price.lean}, l.~224 et 235)
\tierA. Le rapprochement entre les deux dépôts reste textuel \tierC.

\begin{remark}[Contrôles \tierA]
\emph{Contrôle positif.} Avec trois types, $n=(1,2,3)$, $w=(1,2,4)$, $z=(0,1,3)$ et $z'=(1,1,2)$, les deux membres
de l'équation de Price valent $-\frac{13}{36}$ (\lean{ctl_pos_three_types}, l.~249).
\emph{Biais de transmission.} Prenons $n=(1,1)$ et $w=z=(1,2)$, mais une valeur sélective des descendants
$z'=(1,1)$. Alors $\Delta\bar w=-\frac12$, tandis que $\mathrm{Var}/\bar w=\frac16$ : la formule de variance
\emph{échoue}. L'équation de Price, elle, reste vraie, avec $-\frac34=-\frac34$ (\lean{ctl_neg_transmission_bias},
l.~263).
\emph{Le facteur $1/\bar w$ compte.} Sans biais, $\Delta\bar w=\frac16$ mais $\mathrm{Var}(w)=\frac14$
(\lean{ctl_neg_missing_wbar}, l.~280).
\end{remark}

\begin{warnbox}[: un résultat négatif, annoncé comme tel]
Le préenregistrement du programme P7 comportait un item \texttt{fisher\_general} : le théorème fondamental de
Fisher dans la lecture du « changement partiel » de Price (1972), avec une probabilité a priori de $0{,}35$. Cet
item n'a \textbf{pas été atteint}. Seul le cas haploïde à valeurs sélectives constantes est démontré. Ce n'est
pas non plus le cas diploïde multi-allélique, où la monotonie est une inégalité et non cette égalité. Price (1970,
1972) n'a pas été lu ; il est cité d'après LaGatta \cite{ch4:LaG} (p.~69).
\end{warnbox}

% ---------------------------------------------------------------------
\section{Théorie évolutive des graphes : le théorème isotherme}\label{ch4:sec:graphes}\sectionmark{Le théorème isotherme}

\subsection{Une population structurée}
Dans un tissu, une cellule qui se divise ne remplace pas n'importe quelle autre cellule : elle remplace une
\emph{voisine}. La théorie évolutive des graphes décrit cette structure par un graphe orienté pondéré. Les
individus sont les sommets $V=\{1,\dots,n\}$. La \emph{matrice de remplacement} $W$ est à coefficients $\ge0$, et
\emph{stochastique en lignes} : $\sum_jW_{ij}=1$. Le coefficient $W_{ij}$ est la probabilité que le descendant de
$i$ remplace $j$ ; les boucles $W_{ii}>0$ sont permises.

Un \emph{état} est l'ensemble $S\subseteq V$ des sommets mutants. Les mutants ont la valeur sélective $r>0$ et les
résidents la valeur sélective~$1$ ; on note $f_S(i)=r$ si $i\in S$ et $f_S(i)=1$ sinon. Fixons des poids de
sélection $\mu_i>0$, et posons $F(S)=\sum_k\mu_kf_S(k)$. Un pas de la mise à jour \emph{naissance--mort} se
déroule ainsi.
\begin{itemize}[itemsep=1pt]
\item Le sommet $i$ est choisi pour se reproduire avec probabilité $\mu_if_S(i)/F(S)$.
\item Son descendant remplace $j$ avec probabilité $W_{ij}$.
\item Le nouvel état est $S\cup\{j\}$ si $i\in S$, et $S\setminus\{j\}$ sinon.
\end{itemize}
Le cas $\mu\equiv1$ est le processus de Lieberman, Hauert et Nowak \cite{ch4:LHN}. La chaîne a $2^n$ états ; $\emptyset$
et $V$ sont absorbants. On note $P^{+}(S)$ et $P^{-}(S)$ les probabilités que le \emph{nombre} de mutants $|S|$
augmente, resp.\ diminue, d'une unité.

\begin{definition}
La matrice $W$ est \emph{isotherme} si elle est, de plus, stochastique en colonnes : $\sum_iW_{ij}=1$. La
« température » $\sum_iW_{ij}$ de chaque sommet vaut alors~$1$. Elle est \emph{fortement connexe} si, pour tous
$a,b$, il existe un chemin $a=x_0,\dots,x_m=b$, $m\ge1$, avec $W_{x_kx_{k+1}}>0$.
\end{definition}

\begin{leanbox}[: le modèle (\texttt{formal\_cert/Programmes/EvoGraph/BirthDeath.lean}, l.~44--61, 73--74, 90--91)]
\begin{lstlisting}[language=Lean]
noncomputable def fit (r : ℝ) (S : Finset (Fin n)) (i : Fin n) : ℝ := if i ∈ S then r else 1
noncomputable def ind (S : Finset (Fin n)) (i : Fin n) : ℝ := if i ∈ S then 1 else 0
def upd (S : Finset (Fin n)) (i j : Fin n) : Finset (Fin n) := if i ∈ S then insert j S else S.erase j
noncomputable def totFit (μ : Fin n → ℝ) (r : ℝ) (S : Finset (Fin n)) : ℝ := ∑ k, μ k * fit r S k
noncomputable def coef (μ : Fin n → ℝ) (W : Fin n → Fin n → ℝ) (r : ℝ) (S : Finset (Fin n))
    (i j : Fin n) : ℝ :=
  μ i * fit r S i / totFit μ r S * W i j
noncomputable def bdK (μ : Fin n → ℝ) (W : Fin n → Fin n → ℝ) (r : ℝ) (S T : Finset (Fin n)) : ℝ :=
  ∑ i, ∑ j, coef μ W r S i j * (if upd S i j = T then 1 else 0)
noncomputable def upP (μ : Fin n → ℝ) (W : Fin n → Fin n → ℝ) (r : ℝ) (S : Finset (Fin n)) : ℝ :=
  ∑ T, bdK μ W r S T * (if T.card = S.card + 1 then 1 else 0)
def StronglyConnected (W : Fin n → Fin n → ℝ) : Prop :=
  ∀ a b, Relation.TransGen (fun x y => 0 < W x y) a b
\end{lstlisting}
L'auditeur indépendant a jugé que \lean{upP}, \lean{downP} et \lean{bdK} modélisent bien le processus de
Lieberman--Hauert--Nowak. C'est un jugement d'IA, non un audit humain \tierC.
\end{leanbox}

\subsection{Le biais vers l'avant, configuration par configuration}
Pour $S\subseteq V$, on définit les \emph{coupes}
\[
\mathrm{out}(S)=\sum_{i\in S}\sum_{j\notin S}\mu_iW_{ij},\qquad \mathrm{in}(S)=\sum_{i\notin S}\sum_{j\in S}\mu_iW_{ij}.
\]

\begin{lemma}[Les mouvements sont des coupes]\label{ch4:lem:cuts}
On a $P^{+}(S)=r\,\mathrm{out}(S)/F(S)$ et $P^{-}(S)=\mathrm{in}(S)/F(S)$. La chaîne est stochastique, et
$\emptyset$, $V$ sont absorbants. \tierA{} (\lean{up_eq}, \lean{down_eq}, \lean{bd_absorbing},
\texttt{BirthDeath.lean}, l.~148--250).
\end{lemma}

\begin{humanproof}
Les coefficients $c(S,i,j)=\mu_if_S(i)W_{ij}/F(S)$ sont $\ge0$, et leur somme vaut $1$, car $\sum_jW_{ij}=1$ et
$\sum_i\mu_if_S(i)=F(S)$. Quatre cas se présentent selon que $i$ et $j$ sont dans $S$ ou non.
\begin{itemize}[itemsep=1pt]
\item $i\in S$ et $j\notin S$ : un mutant envahit un résident, et $|S|$ augmente. Le poids vaut $\mu_irW_{ij}/F$.
\item $i\notin S$ et $j\in S$ : $|S|$ diminue. Le poids vaut $\mu_iW_{ij}/F$.
\item Dans les deux autres cas, l'état ne change pas.
\end{itemize}
En sommant, on obtient les deux formules. Enfin, $\mathrm{upd}(\emptyset,i,j)=\emptyset$ et $\mathrm{upd}(V,i,j)=V$.
\end{humanproof}

\begin{proposition}[Équilibre des coupes et biais vers l'avant]\label{ch4:prop:bias}
Supposons $\mu$ stationnaire : $\sum_i\mu_iW_{ij}=\mu_j$ pour tout $j$. Alors $\mathrm{out}(S)=\mathrm{in}(S)$ pour
\emph{toute} configuration $S$, et donc
\[
P^{+}(S)=r\,P^{-}(S)\qquad\text{pour tout }S\subseteq V .
\]
En particulier, c'est le cas si $W$ est isotherme et $\mu\equiv1$. \tierA{} (\lean{flow_balance},
\lean{cut_balance}, \lean{forward_bias}, \texttt{BirthDeath.lean}, l.~176--256).
\end{proposition}

\begin{humanproof}
On calcule la même quantité de deux façons. D'une part, $\sum_{i\in S}\sum_{j\in V}\mu_iW_{ij}=\mu(S)$, puisque
$\sum_jW_{ij}=1$. D'autre part, $\sum_{j\in S}\sum_{i\in V}\mu_iW_{ij}=\mu(S)$, par stationnarité. Les deux doubles
sommes contiennent la partie commune $\sum_{i\in S}\sum_{j\in S}\mu_iW_{ij}$. Une fois celle-ci retirée, il reste
$\mathrm{out}(S)=\mathrm{in}(S)$. Le lemme~\ref{ch4:lem:cuts} donne alors $P^{+}=r\,\mathrm{out}/F=r\,\mathrm{in}/F=rP^{-}$.
Si $W$ est isotherme, $\sum_i1\cdot W_{ij}=1=\mu_j$ : c'est la stationnarité de $\mu\equiv1$.
\end{humanproof}

\begin{leanbox}[: le biais vers l'avant (\texttt{BirthDeath.lean}, l.~251--253)]
\begin{lstlisting}[language=Lean]
theorem forward_bias {μ : Fin n → ℝ} {W : Fin n → Fin n → ℝ} (r : ℝ) (hrow : ∀ i, ∑ j, W i j = 1)
    (hstat : ∀ j, ∑ i, μ i * W i j = μ j) (S : Finset (Fin n)) :
    upP μ W r S = r * downP μ W r S := by
\end{lstlisting}
Le biais vers l'avant de Pattni et al. \tierL{}~\cite{ch4:PBRS} (lemme~1, p.~11) et le rapport $p^-/p^+=1/r$ du
théorème~3.2 de Keller--Uğurlu \tierL{}~\cite{ch4:KU} disent la même chose. Remarquons que l'identité s'écrit
sans division : elle reste vraie même si $P^{-}(S)=0$.
\end{leanbox}

\begin{proposition}[Une fonction harmonique sur les configurations]\label{ch4:prop:harmgraph}
Soit $\rho_n(k)=S_k/S_n$ la fonction de fixation de Moran du théorème~\ref{ch4:thm:moran}. Si $\mu$ est
stationnaire, la fonction $h(S)=\rho_n(|S|)$ est harmonique pour la chaîne des configurations en tout
$S\ne\emptyset$. Si de plus une arête $i\in S\to j\notin S$ vérifie $\mu_iW_{ij}>0$, la variance à un pas de $h$ en
$S$ est strictement positive. \tierA{} (\lean{moranFix_card_harmonic}, \lean{moranFix_card_var_pos},
\texttt{BirthDeath.lean}, l.~285--328).
\end{proposition}

\begin{humanproof}
Posons $k=|S|\ge1$. Par l'étape~1 du théorème~\ref{ch4:thm:moran}, $h$ monte de $\Delta^{+}=r^{-k}/S_n$ lors d'un
pas $+1$, et descend de $\Delta^{-}=r^{-(k-1)}/S_n=r\,\Delta^{+}$ lors d'un pas $-1$. Sinon, $h$ ne change pas. Donc
\[
\mathbb E\bigl[h(S')-h(S)\bigr]=P^{+}(S)\Delta^{+}-P^{-}(S)\Delta^{-}=\frac{r^{-k}}{S_n}\bigl(P^{+}(S)-r\,P^{-}(S)\bigr)=0
\]
par la proposition~\ref{ch4:prop:bias}. Pour la variance, le mouvement $i\to j$ a une probabilité $>0$ et change
$h$ de $\Delta^{+}\ne0$.
\end{humanproof}

Ce calcul n'utilise \emph{que} le biais vers l'avant. Il n'affirme jamais que le nombre de mutants est une chaîne
de Markov. Nous verrons au §\ref{ch4:ssec:cycle} que cette affirmation serait fausse.

\begin{lemma}[La connexité force les coupes]\label{ch4:lem:conn}
Supposons $W$ fortement connexe et $\mu>0$. Alors toute partie $S$ non vide et différente de $V$ possède une arête
sortante $i\in S\to j\notin S$ avec $\mu_iW_{ij}>0$. \tierA{} (\lean{cutPositive_of_strongly_connected},
\texttt{BirthDeath.lean}, l.~332).
\end{lemma}

\begin{humanproof}
Prenons $a\in S$, $b\notin S$, et un chemin $a=x_0,\dots,x_m=b$ de poids $>0$. Soit $\ell$ le plus petit indice tel
que $x_{\ell+1}\notin S$. Il existe, car $x_m\notin S$, et l'on a $x_\ell\in S$, car $x_0\in S$. L'arête
$x_\ell\to x_{\ell+1}$ sort de $S$. En Lean, la preuve procède par récurrence sur la dérivation \lean{TransGen}.
\end{humanproof}

\subsection{Le théorème}
\begin{theorem}[Théorème isotherme, mise à jour naissance--mort]\label{ch4:thm:iso}
Soient $n\ge1$, $r>0$, et $W$ isotherme et fortement connexe, avec $\mu\equiv1$. Pour \emph{toute} configuration
initiale $S$,
\[
P^t(S,V)\xrightarrow[t\to\infty]{}\rho_n(|S|)=\frac{1-r^{-|S|}}{1-r^{-n}}\quad(r\ne1),\qquad\frac{|S|}{n}\quad(r=1).
\]
En particulier, un mutant unique placé en \emph{n'importe quel} sommet $v$ se fixe avec la probabilité limite
$\frac{1-r^{-1}}{1-r^{-n}}$, exactement comme dans la population bien mélangée. Plus généralement, la même limite
vaut dès que $\mu W=\mu$, $\mu>0$, et que toute partie propre non vide possède une arête sortante.
\tierA{} (\lean{isothermal_fixation}, \lean{isothermal_fixation_single}, \lean{isothermal_fixation_neutral},
\lean{isothermal_fixation_uniform}, \lean{stationary_fixation}, \texttt{BirthDeath.lean}, l.~356--416 ; \lean{isothermal_fixation_single} est ré-exporté sous le nom \lean{ILENE}).
\end{theorem}

\begin{humanproof}
La proposition~\ref{ch4:prop:harmgraph} fournit une fonction harmonique hors de $\{\emptyset,V\}$, qui vaut $0$ en
$\emptyset$ et $1$ en $V$. Le lemme~\ref{ch4:lem:conn} et la proposition~\ref{ch4:prop:harmgraph} assurent que sa
variance est $>0$ hors des deux états absorbants. Le théorème général d'absorption du §\ref{ch4:sec:wf}
s'applique, sur un ensemble fini d'états quelconque. On l'y transporte en numérotant les $2^n$ configurations de
$0$ à $2^n-1$, avec $\emptyset\mapsto0$ et $V\mapsto2^n-1$.
\end{humanproof}

\begin{leanbox}[: le théorème isotherme (\texttt{BirthDeath.lean}, l.~376--378)]
\begin{lstlisting}[language=Lean]
theorem isothermal_fixation {W : Fin n → Fin n → ℝ} {r : ℝ} (hn : 0 < n) (hr : 0 < r)
    (hW : Isothermal W) (hsc : StronglyConnected W) (S : Finset (Fin n)) :
    Tendsto (fun t => sPow (bdK (fun _ => 1) W r) t S univ) atTop (𝓝 (moranFix n r S.card)) := by
\end{lstlisting}
\emph{Réutilisation.} Le module \texttt{Transport.lean} ne redémontre rien. Le théorème
\lean{two_absorbing_limit} (l.~81) conjugue la chaîne des configurations par une bijection avec
\lean{Fin (M + 1)}, puis applique \emph{tel quel} le théorème \lean{ProgPopGen.fixation_limit} de P7. Les fichiers
\texttt{Absorption.lean} et \texttt{MoranModel.lean} sont réutilisés sans modification.
\emph{Ce qui n'est pas prouvé.} La réciproque n'est pas formalisée : un graphe qui donne la valeur de Moran pour
tout $r$ est-il isotherme ? La proposition~2 de Pattni et al.\ est une équivalence ; seul le sens « suffisant »
est formalisé ici \tierL{}~\cite{ch4:PBRS}. La mise à jour mort--naissance n'est pas formalisée non plus.
\end{leanbox}

\subsection{Pourquoi le nombre de mutants ne se « regroupe » pas}\label{ch4:ssec:cycle}
\begin{center}
\setlength{\unitlength}{1cm}
\begin{picture}(9.5,3.0)
\put(0.5,2.5){\circle{0.5}}\put(0.42,2.42){\small 0}
\put(2.5,2.5){\circle{0.5}}\put(2.42,2.42){\small 1}
\put(2.5,0.5){\circle{0.5}}\put(2.42,0.42){\small 2}
\put(0.5,0.5){\circle{0.5}}\put(0.42,0.42){\small 3}
\put(0.75,2.5){\line(1,0){1.5}}\put(2.5,0.75){\line(0,1){1.5}}
\put(0.75,0.5){\line(1,0){1.5}}\put(0.5,0.75){\line(0,1){1.5}}
\put(3.3,1.4){\parbox{6cm}{\small Le 4-cycle : chaque sommet envoie son descendant sur l'un de ses deux voisins,
avec probabilité $\frac12$. $W$ est doublement stochastique, donc isotherme.}}
\end{picture}
\end{center}

\begin{proposition}[Le 4-cycle \tierA]\label{ch4:prop:cycle}
Sur le 4-cycle, avec $r=2$ et $\mu\equiv1$, les configurations $\{0,1\}$ et $\{0,2\}$ ont toutes deux deux mutants,
mais
\[
P^{+}(\{0,1\})=\tfrac13,\qquad P^{+}(\{0,2\})=\tfrac23,
\]
et le biais vaut $P^{+}=2P^{-}$ dans les deux cas (\lean{cycle4_not_lumpable}, \texttt{Examples.lean},
l.~239--253).
\end{proposition}

\begin{humanproof}
Pour $S=\{0,1\}$, on a $F=2+2+1+1=6$. Les arêtes sortantes sont $0\to3$ et $1\to2$, de poids $\frac12$ chacune,
donc $\mathrm{out}=1$, d'où $P^{+}=2\cdot1/6=\frac13$. De même $\mathrm{in}=1$ et $P^{-}=\frac16$. Pour
$S=\{0,2\}$, les quatre arêtes issues de $0$ et de $2$ sortent toutes, donc $\mathrm{out}=\mathrm{in}=2$,
$P^{+}=\frac46=\frac23$ et $P^{-}=\frac13$.
\end{humanproof}

\begin{leanbox}[: le contre-exemple (\texttt{Examples.lean}, l.~230--231 et 239--243)]
\begin{lstlisting}[language=Lean]
noncomputable def Wcyc : Fin 4 → Fin 4 → ℝ :=
  ![![0, 1 / 2, 0, 1 / 2], ![1 / 2, 0, 1 / 2, 0], ![0, 1 / 2, 0, 1 / 2], ![1 / 2, 0, 1 / 2, 0]]
theorem cycle4_not_lumpable :
    ({0, 1} : Finset (Fin 4)).card = ({0, 2} : Finset (Fin 4)).card ∧
    upP (fun _ => 1) Wcyc 2 {0, 1} = 1 / 3 ∧ upP (fun _ => 1) Wcyc 2 {0, 2} = 2 / 3 ∧
    upP (fun _ => 1) Wcyc 2 {0, 1} = 2 * downP (fun _ => 1) Wcyc 2 {0, 1} ∧
    upP (fun _ => 1) Wcyc 2 {0, 2} = 2 * downP (fun _ => 1) Wcyc 2 {0, 2} := by
\end{lstlisting}
\end{leanbox}

Les probabilités de transition du nombre $|S_t|$ ne sont donc \emph{pas} des fonctions de $|S_t|$. Le nombre de
mutants n'est pas une chaîne de naissance et de mort à taux fixés, au sens du regroupement fort (« strong
lumpability ») de Kemeny et Snell \cite{ch4:KS}. Deux précisions s'imposent \tierC. D'abord, nous ne prétendons pas que le
nombre de mutants soit non markovien pour \emph{toute} loi initiale : seul le regroupement fort est réfuté.
Ensuite, la \emph{chaîne des sauts}, celle où l'on supprime les pas où rien ne se passe, se regroupe bien. Sur le
4-cycle, sachant qu'un mouvement a lieu, les deux configurations montent avec la même probabilité
$P^{+}/(P^{+}+P^{-})=\frac23$, alors que leurs probabilités de rester sur place valent $\frac12$ et $0$ (calcul
rationnel exact \tierX). C'est l'argument de Pattni et al. \tierL{}~\cite{ch4:PBRS}. La preuve Lean ne regroupe
jamais rien : elle travaille configuration par configuration (proposition~\ref{ch4:prop:harmgraph}).

\subsection{Le graphe en étoile : amplifier ou atténuer la sélection}
\begin{center}
\setlength{\unitlength}{1cm}
\begin{picture}(7,1.6)
\put(0.5,0.8){\circle{0.5}}\put(0.42,0.72){\small 1}
\put(3.5,0.8){\circle{0.5}}\put(3.42,0.72){\small 0}
\put(6.5,0.8){\circle{0.5}}\put(6.42,0.72){\small 2}
\put(3.25,0.95){\vector(-1,0){2.5}}\put(1.8,1.1){\small $\frac12$}
\put(0.75,0.65){\vector(1,0){2.5}}\put(1.9,0.25){\small $1$}
\put(3.75,0.95){\vector(1,0){2.5}}\put(4.9,1.1){\small $\frac12$}
\put(6.25,0.65){\vector(-1,0){2.5}}\put(5.0,0.25){\small $1$}
\end{picture}
\[
W_\star=\begin{pmatrix}0&\frac12&\frac12\\[1pt]1&0&0\\1&0&0\end{pmatrix}\qquad
(\text{centre }0,\ \text{feuilles }1,2).
\]
\end{center}

L'étoile est stochastique en lignes, mais sa colonne~$0$ a pour somme $2$ : le centre a la « température » $2$, et
le graphe n'est \emph{pas} isotherme (\lean{star_not_isothermal}) \tierA. Au départ d'un mutant au centre, on a
$P^{+}(\{0\})=P^{-}(\{0\})=\frac12$. Le rapport vaut $1$, et non $r=2$ (\lean{star_bias}) \tierA.

\begin{proposition}[Valeurs exactes sur l'étoile \tierA]\label{ch4:prop:star}
Pour $n=3$, $r=2$ et $\mu\equiv1$, les probabilités limites de fixation valent
\[
\tfrac5{12}\ \text{(mutant au centre)},\qquad \tfrac23\ \text{(mutant sur une feuille)},\qquad
\tfrac7{12}\ \text{(mutant placé uniformément)},
\]
toutes différentes de la valeur de Moran $\rho_3(1)=\frac47$ (\lean{star_fixation_centre},
\lean{star_fixation_leaf}, \lean{star_fixation_uniform}, \lean{star_ne_moran}, \texttt{Examples.lean},
l.~171--196). Avec les poids stationnaires $\mu=(\frac12,\frac14,\frac14)$, on retrouve $\frac47$ au départ du
centre (\lean{starStationary_fixation}, l.~219).
\end{proposition}

La démonstration à la main fait l'objet de l'exercice~\ref{ch4:ex:star}. En Lean, la table des $2^3$ valeurs
$h=(0,\frac5{12},\frac23,\frac23,\frac56,\frac56,\frac{14}{15},1)$ est donnée explicitement (\lean{hStar}, l.~139).
On vérifie qu'elle est harmonique (\lean{star_harmonic}) et de variance positive, puis on la passe dans la
\emph{même} machine d'absorption transportée. Un script indépendant en rationnels exacts redonne la même table
\tierX.

\begin{physbox}[: amplificateurs et atténuateurs]
\tierC{} Comparons à la valeur de Moran $\frac47\approx0{,}571$. Pour un mutant placé uniformément au hasard, l'étoile
\emph{amplifie} la sélection, puisque $\frac7{12}\approx0{,}583>\frac47$. Pour un mutant né au centre, elle
l'\emph{atténue}, avec $\frac5{12}\approx0{,}417$. Il faut lire ces conclusions avec soin. Ce ne sont que
\emph{trois nombres}, pour $n=3$ et $r=2$. Ils ne montrent pas que l'étoile amplifie la sélection pour tout $r$ ou
pour tout $n$. Le préenregistrement qualifiait d'ailleurs l'étoile d'«~amplificateur~» sans réserve. Pour $n=4$,
un script non vérifié donne en rationnels exacts $\frac{392}{1265}$ (centre), $\frac{168}{253}$ (feuille) et une
moyenne $\frac{728}{1265}\ne\frac8{15}$. Ces valeurs relèvent du niveau~\tierX{} et ne sont pas vérifiées par le
noyau. Nous ne citons pas la formule classique de l'étoile, car la source (Lieberman--Hauert--Nowak) n'a pas été
lue.
\end{physbox}

\subsection{La connexité est nécessaire}
\begin{proposition}[Graphe identité \tierA]\label{ch4:prop:identite}
Sur $2$ sommets, la matrice $W=I$ est doublement stochastique, donc isotherme, mais pas fortement connexe. Avec
$r=2$, on a $P^t(\{0\},V)=0$ pour tout $t$, qui ne tend donc pas vers $\rho_2(1)=\frac23$
(\lean{identity_no_fixation}, \texttt{Examples.lean}, l.~289--292).
\end{proposition}
Chaque individu ne remplace que lui-même, et rien ne bouge jamais. L'hypothèse de forte connexité du
théorème~\ref{ch4:thm:iso} ne peut donc pas être supprimée. C'est un \emph{contre-exemple}, et non une réciproque
démontrée \tierC.

\subsection{Une erreur de préenregistrement, racontée honnêtement}
\begin{warnbox}[: ce qui avait été annoncé, et ce qui a été prouvé]
\tierC{} Le 26 octobre 2026, l'item N4 du préenregistrement annonçait : prouver que si le graphe est isotherme,
« \emph{the number of mutants is a birth-death chain with the Moran ratio $r$} », et en déduire la fixation
$(1-r^{-1})/(1-r^{-N})$ en réutilisant P7. Voici ce qui s'est passé.
\begin{itemize}[itemsep=1pt]
\item La \emph{conclusion} est démontrée, et même davantage. Elle vaut pour toute configuration initiale, pour un
mutant en n'importe quel sommet, pour un mutant uniforme, et pour $r=1$ (théorème~\ref{ch4:thm:iso}).
\item La \emph{clause intermédiaire} en italiques est \textbf{fausse} pour la chaîne qui contient les pas où rien
ne se passe. Elle n'a donc pas été formalisée. Le contre-exemple du 4-cycle (proposition~\ref{ch4:prop:cycle}) est
vérifié par le noyau.
\item Ce qui est vrai et prouvé, c'est le biais vers l'avant \emph{configuration par configuration} :
$P^{+}(S)=rP^{-}(S)$ (proposition~\ref{ch4:prop:bias}).
\end{itemize}
Ce n'est pas une découverte. Keller et Uğurlu observent que l'argument de Lieberman--Hauert--Nowak projette la
chaîne sur le nombre de mutants, et que « the projection does not need to be Markovian » \tierL{}~\cite{ch4:KU}
(résumé, p.~4). L'erreur du préenregistrement reproduisait donc une lacune connue d'une preuve publiée. Keller et
Uğurlu écrivent aussi que le théorème « remains unproven till today ». Nous ne reprenons pas cette phrase, car le
lemme~1 de Pattni et al., antérieur, en donne une preuve valide par la chaîne des sauts \tierC. La leçon est de
méthode : énoncer l'objectif, laisser le noyau trancher, puis publier l'écart.
\end{warnbox}

\begin{physbox}[: portée biologique]
\tierC{} Dans ce modèle, une structure équilibrée (doublement stochastique) et fortement connexe ne change pas la
probabilité de fixation d'un mutant de valeur sélective $r$, où qu'il apparaisse. Une structure qui amplifie ou
atténue la sélection doit donc rompre cet équilibre ; l'étoile en est un exemple, pour une seule valeur des
paramètres. Les tissus, où les cellules remplacent leurs voisines, sont la motivation habituelle. Nous ne faisons
cependant aucune affirmation propre à un tissu, et n'utilisons aucune donnée. Le modèle omet la mutation, les
autres règles de mise à jour, les structures qui changent et la dépendance en fréquence.
\end{physbox}

% ---------------------------------------------------------------------
\section{Bilan : ce qui est vérifié, ce qui ne l'est pas}\label{ch4:sec:bilan}\sectionmark{Bilan}

\begin{center}\small
\begin{tabular}{p{0.44\linewidth}p{0.33\linewidth}c}
\toprule
énoncé & déclaration Lean (module) & niveau\\
\midrule
Hardy--Weinberg stationnaire & \lean{hwStep_stationary} (H12) & A\\
identité de l'incrément, monotonie de $\bar w$ & \lean{wbar_increment} (H12) & A\\
$i/N$ harmonique et unique (principe du maximum) & \lean{harmonic_unique} (H11) & A\\
$P^n(i,N)\to h(i)$, chaîne absorbante générale & \lean{fixation_limit} (P7) & A\\
Wright--Fisher : $P^n(i,N)\to i/N$ & \lean{wright_fisher_fixation_limit} (P7) & A\\
décroissance géométrique de $I_n$ & --- (lemme~\ref{ch4:lem:geom}) & C\\
Moran : $P^n(i,N)\to\frac{1-r^{-i}}{1-r^{-N}}$ & \lean{moran_fixation_limit_closed} (P7) & A\\
équation de Price ; $\Delta\bar w=\mathrm{Var}/\bar w$ (haploïde) & \lean{price_equation}, \lean{fisher_haploid_div} (P7) & A\\
théorème de Fisher, lecture « changement partiel » & non atteint & ---\\
biais vers l'avant $P^{+}=rP^{-}$ & \lean{forward_bias} (N4) & A\\
théorème isotherme (sens suffisant) & \lean{isothermal_fixation} (N4) & A\\
4-cycle : $\frac13\ne\frac23$ & \lean{cycle4_not_lumpable} (N4) & A\\
étoile $n=3$ : $\frac5{12}$, $\frac23$, $\frac7{12}$ ; Moran $\frac47$ & \lean{star_fixation_uniform} (N4) & A\\
étoile $n=4$ & script exact & X\\
« probabilité de fixation » $=\lim P^n(i,N)$ & pas d'espace des trajectoires & C\\
lecture biologique (résistance, tissus) & --- & C\\
\bottomrule
\end{tabular}
\end{center}

\paragraph{Rejeu noyau effectué pour ce chapitre.}
Le 2026-10-26, l'agent rédacteur a lancé \texttt{leanchecker} (Lean~4.34.1, chaîne épinglée, deux
cœurs) sur les dix modules cités. Ce sont \lean{Autoresearch.PopGenFisher}, \lean{Autoresearch.PopGenWrightFisher},
\lean{Programmes.PopGen.Absorption}, \lean{WrightFisherFixation}, \lean{MoranModel}, \lean{MoranFixation},
\lean{Price}, ainsi que \lean{Programmes.EvoGraph.Transport}, \lean{BirthDeath} et \lean{Examples}. Les dix
exécutions ont rendu le code~$0$, en $55$ à $131$~s chacune, sans aucune sortie. Les cinq modules jusque-là sans rejeu
(\lean{PopGenFisher}, \lean{PopGenWrightFisher}, \lean{WrightFisherFixation}, \lean{MoranFixation}, \lean{Price}) satisfont donc maintenant la
définition du niveau~A de ce livre. La recherche par \texttt{grep} de \lean{skipKernelTC}, \lean{addDecl},
\lean{native_decide}, \lean{implemented_by}, \lean{sorry}, \lean{admit} et \lean{axiom} dans ces dix sources ne
trouve que deux mentions dans des commentaires de documentation (\texttt{MoranModel.lean}, l.~237 ;
\texttt{Examples.lean}, l.~6). Trois limites sont à noter.
\begin{itemize}[itemsep=1pt]
\item Le rejeu a porté sur les \texttt{.olean} du producteur, et non sur une recompilation propre. Ces fichiers
sont tous datés d'après leur source.
\item Le rejeu n'est pas \texttt{--fresh} : Mathlib et le cœur de Lean ont été chargés, et non rejoués.
\item Le seul contrôle négatif de cette exécution est un module inexistant, refusé avec le code~$1$. Le contrôle
par falsification sous \lean{debug.skipKernelTC}, rejeté par le même binaire, est celui de l'audit N4.
\end{itemize}
Comme le reste du chapitre, cette vérification est l'œuvre d'un agent d'IA ; aucun humain ne l'a faite. La
trace est \texttt{results/kernel\_replay\_2026-10-26/ch4\_popgen\_trace/leancheck\_ch4.json}. Le rejeu de l'ensemble du dépôt, lancé ensuite
(\texttt{results/kernel\_replay\_2026-10-26/replay.jsonl}), a aussi rendu le code~$0$ pour \lean{PopGenFisher}, \lean{PopGenWrightFisher}
et \lean{PopGenImportBoth}.

\paragraph{Ce qui n'est pas formalisé.} Les points suivants restent hors de portée de ce chapitre.
\begin{itemize}[itemsep=1pt]
\item Il n'y a pas de mesure sur l'espace des trajectoires. Une étape naturelle serait d'identifier \lean{kPow}
avec la loi construite par \lean{ProbabilityTheory.Kernel.traj} (Ionescu--Tulcea, dans Mathlib, d'après Marion
\tierL{}~\cite{ch4:Mar}).
\item Ne sont pas formalisés non plus : les temps d'absorption, la limite de diffusion de Kimura, la mutation, la
sélection diploïde multi-allélique, la dépendance en fréquence, la réciproque du théorème isotherme et la mise à
jour mort--naissance.
\item Aucun humain n'a audité l'adéquation des énoncés. Les registres (« ledgers ») de P7 et de N4 signalent
chaque ligne de niveau~A par le drapeau \texttt{LEDGER\_UNAUDITED\_TIER\_A} : il y en a $14$ pour P7 et $11$
pour N4. La porte Elenchus sort avec le code~$1$, et ce n'est pas caché.
\end{itemize}

\paragraph{Pour citer.} Le programme P7 (fixation, Moran, Price) est publié sous doi:10.5281/zenodo.23653850, et
le programme N4 (théorème isotherme) sous doi:10.5281/zenodo.23660071 \cite{ch4:P7,ch4:N4}. Aucun des deux n'a été
relu par des pairs.

% ---------------------------------------------------------------------
\section{Exercices}\label{ch4:sec:exercices}

\begin{exercise}[Hardy--Weinberg]\label{ch4:ex:hw}
Une population a les fréquences génotypiques $(x_{11},x_{12},x_{22})=(0{,}5\,;\,0{,}2\,;\,0{,}3)$.
\begin{enumerate}[label=\alph*),itemsep=0pt]
\item Calculer $p$, puis les fréquences génotypiques après une génération d'accouplement au hasard.
\item Vérifier qu'une deuxième génération ne change plus rien. Quel théorème Lean exprime ce fait ?
\end{enumerate}
\emph{Réponse.} $p=0{,}6$, puis $(0{,}36\,;\,0{,}48\,;\,0{,}16)$. Le théorème est \lean{hwStep_stationary}.
\end{exercise}

\begin{exercise}[Sélection génique]
On prend $w_{11}=(1+s)^2$, $w_{12}=1+s$ et $w_{22}=1$, avec $s>-1$.
\begin{enumerate}[label=\alph*),itemsep=0pt]
\item Montrer que $\bar w(p)=(1+ps)^2$ et que $p'=\dfrac{p(1+s)}{1+ps}$.
\item En déduire que, si $s>0$ et $0<p<1$, la suite des fréquences est strictement croissante et converge vers
$1$.
\end{enumerate}
\emph{Indication.} Pour b), écrire $1-p'=\frac{1-p}{1+ps}$ et itérer l'inégalité $1-p_{t+1}\le\frac{1-p_t}{1+p_0s}$.
\end{exercise}

\begin{exercise}[L'identité de l'incrément]\label{ch4:ex:increment}
\begin{enumerate}[label=\alph*),itemsep=0pt]
\item Démontrer l'identité clé $\bar w+a\,pq=p\,w_{11}+q\,w_{22}$, avec $a=w_{11}-2w_{12}+w_{22}$.
\item En déduire~\eqref{ch4:eq:increment} sans recopier la démonstration du cours.
\item Pour $w=(2,1,1)$ et $p=\frac12$, calculer $\bar w(p)$ et $\bar w(p')$. Vérifier l'identité sur cet exemple.
\item Pour $w=(-1,1,0)$ et $p=\frac12$, montrer que $\bar w$ \emph{diminue}. Quel facteur
de~\eqref{ch4:eq:increment} change de signe ?
\item \emph{Sous-dominance.} Avec $w=(1,\frac12,1)$, montrer que $p=\frac12$ est un équilibre. Calculer $p'$ et
$\bar w(p')$ pour $p=\frac25$, et commenter.
\end{enumerate}
\emph{Réponse.} c) $\bar w=\frac54$ et $\bar w(p')=\frac{34}{25}$. d) $\bar w=\frac14$, $p'=0$ et $\bar w(p')=0$ ;
le facteur $\bar w+pw_{11}+qw_{22}=-\frac14$ est négatif. e) $p'=\frac7{19}$, et $\bar w$ passe de $\frac{19}{25}$
à $\frac{277}{361}$ : la moyenne monte, mais $p$ s'éloigne de l'équilibre.
\end{exercise}

\begin{exercise}[Wright--Fisher à deux gènes]\label{ch4:ex:wf2}
On prend $N=2$.
\begin{enumerate}[label=\alph*),itemsep=0pt]
\item Écrire la matrice $P$ ($3\times3$).
\item Montrer que $P^n(1,1)=2^{-n}$ et que $P^n(1,2)=\frac12\bigl(1-2^{-n}\bigr)$. Retrouver $P^2(1,2)=\frac38$
(\lean{ctl_pos_wf_two_steps}).
\item Comparer la vitesse obtenue avec la borne du lemme~\ref{ch4:lem:geom}.
\end{enumerate}
\emph{Réponse.} c) Ici $\delta=\frac14$, donc la borne vaut $(\frac34)^n$, alors que la vraie décroissance est
$(\frac12)^n$.
\end{exercise}

\begin{exercise}[Wright--Fisher à trois gènes]
On prend $N=3$.
\begin{enumerate}[label=\alph*),itemsep=0pt]
\item Calculer les lignes $1$ et $2$ de $P$. Vérifier que, depuis chaque état intérieur, la probabilité de rester
à l'intérieur vaut $\frac23$.
\item En déduire que $I_n(1)=(\frac23)^n$ exactement, et que $\bigl|P^n(1,3)-\frac13\bigr|\le\frac23\,(\frac23)^n$.
\end{enumerate}
\emph{Réponse.} a) La ligne $1$ est $(\frac8{27},\frac49,\frac29,\frac1{27})$ et la ligne $2$ est
$(\frac1{27},\frac29,\frac49,\frac8{27})$.
\end{exercise}

\begin{exercise}[Moran, $N=3$, $r=2$]\label{ch4:ex:moran3}
\begin{enumerate}[label=\alph*),itemsep=0pt]
\item Calculer $p_1,q_1,p_2,q_2$ et vérifier que $p_i=2q_i$.
\item Résoudre la récurrence de la ruine du joueur, et retrouver $h=(0,\frac47,\frac67,1)$.
\item Montrer que la formule « inversée » $(1-r^i)/(1-r^N)$ donne $\frac17$ en $i=1$, et qu'elle ne vérifie pas
l'équation au premier pas.
\end{enumerate}
\emph{Réponse.} a) $p_1=\frac13$, $q_1=\frac16$, $p_2=\frac4{15}$, $q_2=\frac2{15}$.
\end{exercise}

\begin{exercise}[Grande population]
\begin{enumerate}[label=\alph*),itemsep=0pt]
\item Si $r>1$, montrer que $\rho_N(1)=\frac{1-r^{-1}}{1-r^{-N}}$ décroît vers $1-\frac1r$ quand $N\to\infty$.
\item Si $0<r<1$, montrer que $\rho_N(1)\sim(r^{-1}-1)\,r^N$. Comparer avec la valeur neutre $1/N$.
\item Pour $r=1{,}1$, quelle proportion des mutants avantagés est perdue dans une très grande population ? Pourquoi
cette lecture reste-t-elle de niveau~C ?
\end{enumerate}
\end{exercise}

\begin{exercise}[Équation de Price]
\begin{enumerate}[label=\alph*),itemsep=0pt]
\item On prend $n=(1,2,3)$, $w=(1,2,4)$, $z=(0,1,3)$ et $z'=(1,1,2)$. Calculer $\bar w$, $\bar z$, $\bar z'$,
$\mathrm{Cov}_n(w,z)$ et $\mathbb E_n[w\Delta z]$, et vérifier que les deux membres valent $-\frac{13}{36}$.
\item Pour deux types de fréquences $p$, $1-p$, montrer que $\mathrm{Var}_n(w)=p(1-p)(w_1-w_2)^2$.
\item On prend $n=(1,1)$, $w=z=(1,2)$ et $z'=(1,1)$. Montrer que $\Delta\bar w\ne\mathrm{Var}_n(w)/\bar w$, alors que
l'équation de Price reste vraie. Quelle hypothèse du corollaire~\ref{ch4:cor:haploide} manque ?
\end{enumerate}
\emph{Réponse.} a) $\bar w=\frac{17}6$, $\bar z=\frac{11}6$ et $\bar z'=\frac{29}{17}$, donc
$\bar z'-\bar z=-\frac{13}{102}$ ; b) se vérifie par un calcul
direct ; c) $\Delta\bar w=-\frac12$ et $\mathrm{Var}/\bar w=\frac16$ ; il manque l'absence de biais de
transmission.
\end{exercise}

\begin{exercise}[Le graphe complet est le modèle de Moran]
On prend $W_{ij}=\frac1n$ pour tous $i,j$ (boucles comprises) et $\mu\equiv1$.
\begin{enumerate}[label=\alph*),itemsep=0pt]
\item Vérifier que $W$ est isotherme et fortement connexe.
\item Pour $|S|=k$, calculer $\mathrm{out}(S)$, $\mathrm{in}(S)$ et $F(S)$. En déduire que $P^{+}(S)$ et $P^{-}(S)$
sont exactement les taux $p_k$, $q_k$ du modèle de Moran (\lean{wellMixed_up}, \lean{wellMixed_down},
\texttt{Examples.lean}, l.~62--81).
\item En déduire que, ici, la chaîne des configurations se regroupe exactement sur la chaîne de Moran.
\end{enumerate}
\emph{Réponse.} b) $\mathrm{out}=\mathrm{in}=k(n-k)/n$ et $F=rk+n-k$.
\end{exercise}

\begin{exercise}[Le 4-cycle, à la main]
\begin{enumerate}[label=\alph*),itemsep=0pt]
\item Vérifier que la matrice du 4-cycle est doublement stochastique.
\item Retrouver $P^{\pm}(\{0,1\})=(\frac13,\frac16)$ et $P^{\pm}(\{0,2\})=(\frac23,\frac13)$.
\item Calculer, pour chacune des deux configurations, la probabilité de monter \emph{sachant qu'un mouvement a
lieu}, et commenter.
\item Que donne le théorème isotherme pour un mutant unique, quand $r=2$ ?
\end{enumerate}
\emph{Réponse.} c) $\frac23$ dans les deux cas. d) $\rho_4(1)=\frac{1-1/2}{1-1/16}=\frac8{15}$.
\end{exercise}

\begin{exercise}[Probabilité de fixation sur l'étoile, à la main]\label{ch4:ex:star}
On prend l'étoile à trois sommets (centre $0$, feuilles $1$ et $2$), avec $W_{0j}=\frac12$, $W_{j0}=1$, $r=2$ et
$\mu\equiv1$. Par symétrie, on pose $x_0=h(\{0\})$, $x_1=h(\{1\})=h(\{2\})$, $y=h(\{0,1\})=h(\{0,2\})$ et
$z=h(\{1,2\})$.
\begin{enumerate}[label=\alph*),itemsep=0pt]
\item Depuis $\{0\}$ : montrer que $F=4$, que l'on passe à $\{0,j\}$ avec probabilité $\frac12$ (au total) et à
$\emptyset$ avec probabilité $\frac12$. En déduire que $x_0=\frac y2$.
\item Depuis $\{1\}$ : montrer que l'on passe à $\{0,1\}$ avec probabilité $\frac12$, à $\emptyset$ avec
probabilité $\frac18$, et que l'on reste sur place avec probabilité $\frac38$. En déduire que $x_1=\frac45y$.
\item Depuis $\{0,1\}$ : montrer que l'on passe à $V$ avec probabilité $\frac15$ et à $\{1\}$ avec probabilité
$\frac15$, et en déduire que $2y=1+x_1$. Depuis $\{1,2\}$ : montrer que $z=\frac45+\frac{x_1}5$.
\item Résoudre le système, et retrouver $x_0=\frac5{12}$, $x_1=\frac23$, $y=\frac56$ et $z=\frac{14}{15}$, puis la
valeur $\frac7{12}$ pour un mutant uniforme. Comparer à $\frac47$.
\item Reprendre a) avec $\mu=(\frac12,\frac14,\frac14)$. Vérifier que $\mu W_\star=\mu$, puis que $x_0=\frac47$
(\lean{starStationary_fixation}).
\end{enumerate}
\emph{Réponse.} d) On obtient $y=\frac12+\frac25y$, d'où $y=\frac56$.
\end{exercise}

\begin{exercise}[Lire un énoncé Lean]
\begin{enumerate}[label=\alph*),itemsep=0pt]
\item On prend $N=2$. Montrer que les noyaux absorbants pour lesquels $j\mapsto j/2$ est harmonique sont
exactement ceux de la forme $P(1,0)=P(1,2)=a$, $P(1,1)=1-2a$, avec $a\in[0,\frac12]$. Montrer que l'hypothèse
\lean{hvar} de \lean{fixation_limit} équivaut à $a>0$, et que dans ce cas $P^n(1,2)=\frac12\bigl(1-(1-2a)^n\bigr)$.
Que se passe-t-il pour $a=0$ ?
\item Pourquoi un fichier qui contient \lean{sorry} compile-t-il avec le code de retour $0$ ? Quelles sont les trois
vérifications qui constituent la porte du niveau~A dans ce livre ?
\item L'énoncé \lean{isothermal_fixation} parle-t-il de probabilités ? Rédiger en une phrase le pont \tierC{} qui
le relie à la « probabilité de fixation ».
\end{enumerate}
\emph{Indication.} a) La variance à un pas en $1$ vaut $2a\cdot\frac14=\frac a2$. Pour $a=0$, on retrouve le
noyau identité, c'est-à-dire le contrôle négatif \lean{ctl_neg_identity_no_absorption}.
\end{exercise}

\label{ch4:fin}
\fussy
